% ssmPart2meps.tex - for submission to MEPS. 18th June 2019.
% ssmPart2.tex - second manuscript. 22nd October 2015.

% ssmDraft.Snw - doing draft as Sweave (but importing R results. 9th Sept 2014.
% ssmDraft.tex - size-spectra methods manuscript, working draft.
%  Andy using an earlier submission to Methods in Ecology and Evolution
%  as a template (which came from my Journal of Animal Ecology paper),
%   since we may submit there. Sept 2014.
%  20th May 2014

%%%  FOR TYPESETTERS
%%%  This is my latex file for the manuscript. Comments start with %. Please
%%%  contact me with any queries, especially if you are going to run Latex
%%%  rather than just use the text from this file. It should run okay,
%%%  though you will need the packages epsfig, rotating and possibly natbib.
%%%   Andrew.Edwards@dfo-mpo.gc.ca
%%%
%%%  To make the .pdf file, I run:
%%%  latex ssmDraft.tex
%%%  dvips ssmDraft
%%%  ps2pdf ssmDraft.ps ssmDraft.pdf

\documentclass[12pt, letter]{article}
\textheight 213mm
\topmargin -10mm
\addtolength{\textwidth}{1.0in}
\addtolength{\oddsidemargin}{-0.5in}
% \usepackage{Sweave}
\usepackage{epsfig}
\usepackage{rotating}           % For sideways table
\usepackage[pagewise]{lineno}
\usepackage{amsmath}       % for \text for x_min, for \dfrac
\usepackage{cancel}        % for \cancel
\usepackage{natbib}
\usepackage{array}         % for fixed with tables
\usepackage[table]{xcolor} % apparently best to load last. For Appendix tables.
\usepackage{forloop}

\usepackage{mathptmx}      % Bolder Times plus has math fonts (as for hake)

\usepackage{minitoc}       % See
      % https://tex.stackexchange.com/questions/419249/table-of-contents-only-for-the-appendix
% Make the "Part I" text invisible [ignore main text in TOC]
\renewcommand \thepart{}
\renewcommand \partname{}

\usepackage{tocloft}       % For ToC
\addtolength{\cftsubsecnumwidth}{6pt} % For spacing after section numbers in
                                    %  Supp Info TOC.
% \makeatletter \renewcommand{\@dotsep}{4.5} \makeatother % for dots, didn't work.

\bibliographystyle{natbib}


\linenumbers


% \pagestyle{headings}
\newcommand{\eb}{\begin{eqnarray}}
\newcommand{\ee}{\end{eqnarray}}
\newcommand{\xmin}{x_{\mathrm{min}}}
\newcommand{\xmax}{x_{\mathrm{max}}}
\newcommand{\logten}{\log_{\mathrm{10}}}
\newcommand{\logtwo}{\log_{\mathrm{2}}}
\newcolumntype{L}[1]{>{\raggedright\let\newline\\\arraybackslash\hspace{0pt}}p{#1}}    % From http://tex.stackexchange.com/questions/12703/how-to-create-fixed-width-table-columns-with-text-raggedright-centered-raggedlef   for fixed with columns     though I changed m{#1} to p{#1}

\newcommand{\dx}{\mbox{d}x}
% \newcommand{\standHist}{\ref{fig:dataCode/standHist}}
% \newcommand{\fit}{\ref{fig:dataCode/fitting2a}}
\newcommand{\fitMLEhists}{\ref{fig:binFitting/fitMLEhists}}
\newcommand{\fitMLEconfs}{\ref{fig:binFitting/fitMLEconfs}}
\newcommand{\fitMLEhistsXminSixteen}{\ref{fig:binFitting/xmin16/fitMLEhistsXmin16}}
\newcommand{\fitMLEconfsXminSixteen}{\ref{fig:binFitting/xmin16/fitMLEconfsXmin16}}
\newcommand{\fitMLEhistsCutOffSixteen}{\ref{fig:binFitting/xCutOff16/fitMLEhistsCutOff16}}
\newcommand{\fitMLEconfsCutOffSixteen}{\ref{fig:binFitting/xCutOff16/fitMLEconfsCutOff16}}
\newcommand{\MEE}{erpbb17}     % bib citation for Edwards et al. (2017) MEE paper
\newcommand{\wordcount}{\input{\jobname.wc}}  % To automate wordcount

% \def\fracd#2{\frac{\displaystyle #1}{\displaystyle #2}}

% Had this for Ecology, probably tweak for JAE:
% \bibpunct{(}{)}{,}{a}{}{,}   % see natbib.pdf, p13

\newcommand\onefig[2]{    % filename is #1, text is #2
  \begin{figure}[tp]
  \begin{center}
   % \includegraphics[width=6in,height=7in,keepaspectratio=TRUE]{#1.eps} \\  % RH much better control
  \epsfxsize=6in
  \epsfbox{../../fitting-size-spectra/codeBinStructure/#1.eps}
  \end{center}
  \caption{#2 }
  \label{fig:#1}
  \end{figure}
  % \clearpage
}

\newcommand\onefigshift[2]{    % filename is #1, text is #2
  \begin{figure}[tp]
  \begin{center}
   % \includegraphics[width=6in,height=7in,keepaspectratio=TRUE]{#1.eps} \\  % RH much better control
  \epsfxsize=6in
  \epsfbox{../../fitting-size-spectra/codeBinStructure/#1.eps}
  \end{center}
  \vspace{-10mm}
  \caption{#2}
  \label{fig:#1}
  \end{figure}
  % \clearpage
}

\newcommand\isdfig[2]{    % filename is #1, text is #2
  \begin{figure}[tp]
  \begin{center}
  \vspace{-5mm}
   % \includegraphics[width=6in,height=7in,keepaspectratio=TRUE]{#1.eps} \\  % RH much better control
  % \epsfxsize=6.5in
  \epsfbox{../../fitting-size-spectra/codeBinStructure/nSeaFung/IBTS-ISDfigs/IBTS-ISD#1.eps}
  \end{center}
  \vspace{-5mm}
  \caption{#2}
  \label{fig:ISD-#1}
  \end{figure}
  \clearpage
}


\newcommand\simpfig[2]{    % filename is #1, text is #2. No epsfxsize
  \begin{figure}[tp]
  \begin{center}
   % \includegraphics[width=6in,height=7in,keepaspectratio=TRUE]{#1.eps} \\  % RH much better control
  % \epsfxsize=6in
  \epsfbox{#1.eps}
  \end{center}
  \caption{#2 }
  \label{fig:#1}
  \end{figure}
  \clearpage
}


\newcommand\twofig[4]{   % figure #1 under #2, caption text #3, width #4
  \begin{figure}[tp]     %  label will be #1
  \centering
  \epsfxsize=#4 in
%  \epsfysize=3.5in
  \begin{tabular}{c}
  %	\includegraphics[width=6in,height=3.5in,keepaspectratio=TRUE]{#1.eps} \\  % RH much better control
  %	\includegraphics[width=6in,height=3.5in,keepaspectratio=TRUE]{#2.eps}
  \epsfbox{#1.eps}
  \vspace{-15mm}\\
  \epsfbox{#2.eps}
  \end{tabular}
  \caption{#3}
  \label{fig:#1}
  \end{figure}
  \clearpage
}

\newcommand\threefig[4]{    % figure #1 then #2 then #3,
  \begin{figure}[htp]       %  caption text #4, label will be #1
  \centering
  \begin{tabular}{c}
%	\includegraphics[width=6in,height=2.5in,keepaspectratio=TRUE]{#1.eps} \\  % RH much better control
%	\includegraphics[width=6in,height=2.5in,keepaspectratio=TRUE]{#2.eps} \\
%	\includegraphics[width=6in,height=2.5in,keepaspectratio=TRUE]{#3.eps}
  \vspace{-20mm}
  \epsfbox{#1.eps} \\
  \vspace{-20mm}
  \epsfbox{#2.eps} \\
   \vspace{-20mm}
  \epsfbox{#3.eps}
  \end{tabular}
  \caption{#4}
  \label{fig:#1}
  \end{figure}
}


\renewcommand{\baselinestretch}{1.6}
\renewcommand{\refname}{Literature Cited}
%TC:macroword \onefig 1
%TC:macroword \onefigshift 3
%TC:macroword \simpfig 3
%TC:macroword \simpfigshift 3
%TC:macroword \twofig 3



\begin{document}

%TC:ignore

% Most useful options (with defaults):
% echo        = TRUE     - includes R code in output file
% keep.source = FALSE    - when echoing, if TRUE then original source is copied to the file, otherwise deparsed source is echoed.
% eval        = TRUE     - if FALSE then chunk is not evaluated
% results     = VERBATIM - R output included verbatim, if TEX output is already proper latex and included as is,
%                          if HIDE then all output is completely suppressed (but the code executed - good for admb) results options should all be lower case (else get warnings)
% pdf         = TRUE     - whether .pdf figures shall be generated
% eps         = TRUE     - whether .eps figures shall be generated
% strip.white = FALSE    - if true then blank lines at beginning and end of output are removed. If all, then all blank lines are removed.
% width       = 6        - width of figures in inches
% height      = 6        - height of figures in inches
% fig         = FALSE    - whether the code chunk produces graphical output (only one per chunk)

% \setkeys{Gin}{width=6in}     % from googling sweave figure bigger.
%  It will set this for the rest of document [doesn't width do that in the above?]


% From MEE website, 7/4/15:

% STANDARD ARTICLES
% Standard papers should describe new methods and how they may be used. We place emphasis on methods that are applicable as broadly as possible.
% Structure:
% Our maximum recommended word count for standard papers is 6000-7000 words (including tables, figure captions and references).

\doparttoc % Tell to minitoc to generate a toc for the parts
\faketableofcontents % Run a fake tableofcontents command for the partocs

% \part{} % Start the document part
% \parttoc % Insert the document TOC


\noindent {\Large \bf Accounting for the bin structure of data when fitting
  size spectra}

\vspace{7mm}

\noindent {\large Andrew M.~Edwards$^{1,2,\star}$, James P. W. Robinson$^{2,3}$, Julia
  L. Blanchard$^{4}$, \\Julia K. Baum$^{2}$ and Michael J. Plank$^{5,6}$}

\medskip

{\it
\noindent $^1$Pacific Biological Station, Fisheries and Oceans Canada,
  Nanaimo, British Columbia, V9T 6N7, Canada. %  3190 Hammond Bay Road,
\noindent $^2$Department
  of Biology, University of Victoria, Victoria, British
  Columbia, V8W 2Y2, Canada. % PO Box 1700 STN CSC,
\noindent $^3$Lancaster Environment Centre, Lancaster University, Lancaster,
  LA1 4YW, U.K.
\noindent $^4$Institute for Marine and Antarctic Studies,
  University of Tasmania, Hobart, TAS 7001, Australia. % Private Bag 129,
\noindent $^5$School of Mathematics and Statistics, University of Canterbury,
  Christchurch, New Zealand.
\noindent $^6$Te P\={u}naha Matatini, a New Zealand Centre of
  Research Excellence, University of Auckland, Auckland 1011, New Zealand.}

\medskip

\noindent $^\star$Corresponding author email: Andrew.Edwards@dfo-mpo.gc.ca


\medskip

%\noindent Correspondence author: Andrew M.~Edwards, Pacific Biological Station,
%Fisheries and Oceans Canada, 3190 Hammond Bay Road, Nanaimo, British Columbia,
%V9T 6N7, Canada.\\
%Tel: +1 250 756 7146. Fax: +1 250 756 7053. Email:
%Andrew.Edwards@dfo-mpo.gc.ca

\medskip

\noindent Running page header: Size spectra for binned data % 3-6 words

\medskip

\section*{Abstract}

Size spectra are often recommended as tools for detecting the
response of marine communities to fishing or to management measures.
A size spectrum succinctly describe how a
property, such as abundance or biomass, varies with body size in an ecological
community.
Required
data are often collected in binned form, such as numbers of individuals
in 1-cm length bins.
Numerous methods have been employed to fit size spectra, but most give
inaccurate estimates when tested on simulated data, and none account for the
data's bin structure (breakpoints of bins).
Here, we use eight methods to fit an annual size-spectrum
exponent, $b$, to an example data set (30~years of data from the North Sea
International Bottom Trawl Survey).
The methods give conflicting conclusions regarding $b$ declining
(the size spectrum steepening) through time, and so any resulting advice
to ecosystem managers will be highly dependent upon the method used.
Using simulated data we show that ignoring the bin structure gives inaccurate
estimates of $b$,
even for high-resolution data. However, our extended likelihood method, that
explicitly accounts for the bin structure, accurately estimates $b$ and its
confidence intervals, even for coarsely collected data.
We develop a novel visualisation method that accounts for the bin
structure and associated uncertainty, give
recommendations concerning different data types, and
provide documented R code to encourage use of our methods.
This work is also relevant to wider applications where
a power-law distribution (the underlying distribution for a size spectrum) is
fitted to binned data.

% 250 words (max is 250)

%We then investigate how not accounting for the bin structure
%influences estimates of $b$.
%Our methods have application wherever
% in other ecological contexts then accounting for the bin
%structure may be necessary.


% \noindent [196 words, maximum is 200]
% Only the likelihood method can be extended to deal with binned data.


%in particular the conversion of length data to body-mass data using
%species-specific length-weight relationships.
%However, all ignore the `bin structure' of the data -- the data are
%numbers of individual fish in 1-cm length bins, that get
%converted to individual body masses using

% Thus it may be better (and possibly easier and cheaper) to collect coarsely binned data whose bin structure is properly accounted for, rather than high-resolution data whose bin structure is overlooked.

\vspace{10mm}
% 3-8 key words, in order
\noindent {\it Key words:} abundance size spectrum, biomass size spectrum,
individual size distribution, ecosystem-based fisheries management,
ecosystem indicators, truncated Pareto distribution, power-law distribution

% MEE says maximum of 6000-7000 words
\medskip
\noindent Word count (including references, tables and figure legends): \wordcount
% this is to automate the word count

% Obtain with:
% texcount -1 -sum -relaxed ssmPart2meps.tex > ssmPart2meps.wc


\medskip

\noindent For submission to \emph{Marine Ecology Progress Series}

\noindent \today

%TC:endignore

\clearpage

\section{Introduction}

Ecosystem-based fisheries management requires general indicators to accurately
assess ecosystem states across space and time \citep{jenn05, srjfg05}.
Size spectra quantify how a property (such as abundance or biomass) varies with
body size in a community \citep{rg96}.
They are readily fitted to simple body-size data, and have
therefore been widely applied to aquatic ecosystems
(e.g.~\citealt{dpmg04, dgpr05, bblhw12}).

%One indicator is the size spectrum, which
%quantifies how a property (such as abundance or biomass) varies with body size
%in a community \citep{rg96}. It is applicable to different community types and
%is readily estimated using simple body-size data.

% From Julia Ba.:

%Size spectra quantify how a property
%(such as abundance or biomass) varies with body size in a community (Rice and
%Gislason, 1996). They are readily estimated using simple body-size data, and have
%therefore been widely implemented in aquatic ecosystems (INSERT several
%refs. showing broad usage e.g. James can help - one from lakes, one from
%temperate marine, James' from coral reefs).
%JLB's:
%Empirical and theoretical analyses of fish community size spectra have led to the development of different types of size-based indicators (such as slope, intercept, shape, size diversity and proportion of large fish) to assess the ecosystem effects of fishing.

%The slope of the size spectrum has received most attention due to its relative robustness compared to other size-based indicators under different types of fisheries management scenarios (Piet and Jennings, 2005).

The slope of the size spectrum is a
recommended ecosystem indicator because of its apparent responsiveness to
fishing impacts.
\citet{pj05} tested the response to fishing of eight potential fish community
indicators and found that only the slope of the biomass size spectrum reliably
detected the
effects of spatial management measures (partial closure of an area of
the North Sea). \citet{jd05} similarly recommended the use of size spectra to
provide
surveillance of the status of fish communities. They pointed out difficulties in
justifying reference points (targets to be aimed for and limits to be avoided),
and instead recommended reference directions that monitor trends in
size-spectrum slopes through time. Such trends are commonly investigated in
size-spectrum studies (e.g.~\citealt{blan05, bblhw12}).
Using 188 plausible multispecies size-structured fish community
models, \citet{tllcj15} found that the slope of the abundance size spectrum was
the most responsive indicator of the four indicators tested.
Size-spectra parameters
have been selected by the European Union as one of several indicators used to
evaluate sea-floor integrity \citep{europe10, rice12}, and size-spectrum models
continue to be widely applied
(e.g.~\citealt{tfp14, amln16, ssmarp16, wsj18, zcxxr18, mnwb18, va18, zhou19}).
Such applications require
consistent calculation methods, especially when used for providing practical
advice to fisheries (or ecosystem) managers.
However,
%despite the popularity and broad
% implementation of size spectra,
most commonly employed methods for fitting
size-spectra have been shown to give inaccurate estimates for simulated data
\citep{\MEE}, underscoring the need for close evaluation of methods.


% ... Body size has the potential to become a `first principle of assessment,
%management and conservation of ecosystem status based on formal mathematical
%heory' \citep{pb10}.

Data used in size-spectra studies are often only collected in binned form,
whereby the
`bin structure' (breakpoints of the bins, \citealt{monn16}) is determined during
data collection. For example, \citet{dgpr05} analysed fish lengths collected in
bins of width 0.5~cm, 1~cm and 5~cm. For benthic invertebrates in the
North Sea,
\citet{mj06} sorted individuals into sixteen $\logtwo$ body-mass bins, resulting
in counts and total biomass for each bin.
Similarly, in underwater visual surveys, fish lengths might be recorded to
the nearest 1~cm \citep{mcco16} or 5~cm \citep{rwnz11},
generating size-frequency data of counts per length bin.
For historical data, all that may be available is
digitisation of points from a
published figure and the bin structure has to be
algebraically determined (e.g.~\citealt{edwa11} in a related context).

Alternatively, data are sometimes binned during analysis. For example,
\citet{dpmg04} and \citet{bblhw12} binned data into 5-cm length bins to
calculate size spectra for fish communities in Fiji and the eastern Bering Sea,
respectively. And some size-spectra calculation methods require binning of
data to a
prescribed resolution \citep{\MEE}.
% In the related fisheries context of
% age-structured stock-assessment models, \citet{monn16} found that wider length
% bins can substantially affect estimates of growth parameters but not estimates
% of management quantities.

Here we investigate how the bin structure of data affects the fitting of
size spectra, and whether the bin structure needs to be
explicitly accounted for.
To motivate our work and understand issues that can arise when fitting size
spectra to a complex binned data set,
we first apply eight existing methods to fit annual size spectra to 30 years of data from the
North Sea International Bottom Trawl Survey (IBTS).
%The data are counts per hour of trawling of individuals (identified by species) in 0.5~cm or 1~cm length bins.
The data are the numbers of individuals of each species within each length bin
that were caught per hour of trawling.
% The data are numbers caught per hour of trawling of each species within each
% length bin.
We use the IBTS data set as an example data set for several reasons: (i) it is an open-access
data set downloadable from the International Council for the Exploration of the
Sea (ICES) website, (ii) it is an extensive quality-controlled data set
consisting of many species across multiple years,
(iii) we could use the data-processing protocols that were fully
described by \citet{ffrr12} for their calculation of the Large Fish Indicator,
(iv) we could use the species-specific length-weight coefficients collated
by \citet{ffrr12}.

For each year of data we estimate the size-spectrum exponent, $b$. Explicitly
this is the exponent of the individual size distribution (ISD) and
is related to the slope of the size spectrum \citep{weke07, \MEE}. We find
that five methods estimate a significant declining trend in $b$, corresponding
to a steepening size spectrum (as can be expected from fishing pressure,
e.g.~\citealt{rg96} and \citealt{dgpr05}). The remaining three methods
find no significant trend in $b$. Thus, the ecological conclusions differ depending on the
method used.

However, none of the methods properly take into account the bin structure of the
data.
They implicitly assign a single length to all counts within
a length bin, rather than acknowledging that counts can represent a range of
values within the length bin. The binned lengths are subsequently converted to
body masses through species-specific length-weight relationships.
We illustrate the problems that occur, which are compounded due to
most of the methods requiring further binning to estimate $b$.
This motivates investigation of methods to properly account for the bin
structure.

The likelihood method is the only method that can be extended to properly deal
with binned data. We therefore develop and test such an extended
method, which deals with (i) the bin structure of length or body-mass data, and
(ii) situations giving non-integer counts of fish. The latter occurs because the
counts (for the IBTS data)
are standardised to represent the number of fish caught per hour of
trawling, and also arises with other types of data (e.g.~\citealt{rwnz11}).
Using simulated body-mass data we find that applying likelihood
without taking into account the bin structure of the data yields inaccurate
estimates of $b$, even for high-resolution data. However, the extended
likelihood method that takes into account the bin
structure is accurate even
for low-resolution data. Hence it may not be worth the effort and expense of
collecting high-resolution data if the bin structure is not accounted for.

We derive the likelihood
function that fully accounts for species-specific length-weight
relationships. This allows the bin structure for length data to be properly
accounted for when converting lengths to weights.
Applying this function to the IBTS data, we find consistently greater
(less negative) estimates of $b$ than for the original likelihood method.
Finally, we develop a new method for plotting such data and the resulting
ISD that explicitly accounts for the
bin structure and associated uncertainty.

Our results strengthen and extend those from \citet{\MEE}, where we recommended the use of
maximum likelihood estimation over the seven other previous methods based on
simulated (rather than real) data. Here, only the likelihood method can be
properly extended to deal with the uncertainty associated with binned data.
% Mike's:
% Here, we develop a new method for fitting size-spectra that explicitly accounts
% for the uncertainty associated with binning and/or measurement data. We show
% that this method is more accurate that existing methods and .
% Andy: I included the uncertainty bit, but including all of it gets a bit
% repetivie I think.

We finish with recommendations concerning which version of the likelihood method
to use for particular types of data. To encourage and enable others to implement
our methods, we provide fully documented and usable R code (as Supplementary
Material and currently at

\noindent {\tt
  https://github.com/andrew-edwards/fitting-size-spectra/tree/manuscript2}).


\section{Materials and Methods}

\subsection{Demersal fish data}

For our example data set we
we obtained data collected by the North Sea IBTS from the ICES online Database
of Trawl Surveys ({\tt
  http://www.ices.dk/marine-data/data-portals/Pages/DATRAS.aspx}).
We downloaded the number of
organisms caught per hour for each species and
length class for all surveys in Areas~1 to~7 carried out in Quarter~1
from 1986-2015 (see Supplementary Material for full details).
All surveys followed standardized
gear- and data-collection protocols \citep{ices15}.
We followed
\citet{ffrr12} in removing all non-demersal fish.
During surveys, lengths were recorded as rounded
down to the nearest 1~cm [or 0.5~cm for Atlantic Herring (\emph{Clupea
  harengus}) and European Sprat (\emph{Sprattus sprattus})].
We averaged counts of the same year/species/length-class combination
across the seven areas.
For simplicity this ignores potential differences in
sampling effort due to variations in haul duration, trawl speed and net area,
and fits a common size spectrum for the whole region.

For each species, $s$, the length-weight relationship is
\eb
w = \alpha_s l^{\beta_s},
\label{lw}
\ee
where $w$ is the estimated body mass (g) of an individual, $l$ is the known
body length (cm), and $\alpha_s$ and $\beta_s$ are species-specific parameters
published by \citet{ffrr12}. Following \citet{ffrr12}, for $l$ we initially use
the `Length class' value (see
Table~\ref{tab:egData} for examples), which represents the minimum of possible
lengths for a given fish (i.e.~a `Length class' of 35~cm represents fish in the
range 35-36~cm). This is consistent with, for example, \citet{blan05} and
\citet{dgpr05}, the latter also doing this for 5-cm and 10-cm length
bins. For example, a 45-cm Smallspotted Catshark is assigned a calculated
body mass of $\alpha_s l^ {\beta_s} = 0.0031 \times 45^{3.0290} = 315.46$~g
(first row of Table~\ref{tab:egData}). In 1986 there were 0.00714~h$^{-1}$ of such fish caught,
resulting in a total biomass per hour of trawling of
2.25~g~h$^{-1}$ (Table~\ref{tab:egData}).

Following calculations of size spectra for earlier IBTS data \citep{pj05} and
Celtic Sea groundfish surveys \citep{blan05}, we remove all resulting body
masses estimated to be $<4$~g because small fish will not be fully selected by
the gear.  The resulting dataset (Table~\ref{tab:egData}) contains the required
information to estimate the size-spectrum exponent of the community for each
year. It consists of a dataframe with 42,298 rows, where each row is a unique
combination of year, species and length class. Each year has between 57 and 81
unique species.


% **I did have: `In other circumstances, length classes may be at 5-cm resolution, in which case it is likely that midpoints would be used [Julia Bl. - do you have a reference, or can we say personal observation?]'. Be good to have a ref for using midpoints. Keep a look out.

\subsection{Individual size distribution}

We designate the random variable $X$ to represent the mass (g) of an individual fish
or other organism. Then, considering $X$ to come from a bounded power-law (PLB)
distribution, the probability density function for $X$ is \eb f(x) = C
x^b,~~~~\xmin \leq x \leq \xmax
\label{PLB}
\ee where \eb C = \left\{ \begin{array}{ll} \dfrac{b + 1}{\xmax^{~b + 1} -
    \xmin^{~b + 1}}, \vspace{2mm} & b \neq -1, \\ \dfrac{1}{\log \xmax - \log
    \xmin}, & b = -1,
  \end{array}
  \right.
\label{PLB.C}
\ee $x$ represents possible values of $X$, $\log$ is the natural logarithm, $b$
is the size-spectrum exponent, and $\xmin$ and $\xmax$ are the minimum and
maximum possible values of the data with $0 < \xmin < \xmax$ \citep{\MEE}. The
exponent $b$ is expected to be negative, and is the quantity of interest when
analysing size spectra. Equation~(\ref{PLB}) is the ISD,
and the resulting biomass size spectrum (biomass density function) for a
community of $n$ individuals is
\eb
B(x) = n C x^{b+1},~~~~\xmin \leq x \leq \xmax.
\label{biomass}
\ee
For length size spectra $x$ would represent lengths, but (\ref{biomass}) cannot
be directly used \citep{\MEE}.

\subsection{Estimating $b$ for each year of the demersal fish data}

We apply eight methods for fitting size spectra that have been
previously used in papers and in the
R package {\tt mizer} \citep{sba14}. The methods were
documented by \citet{\MEE}. Here, the Supplementary Material explains how
the methods need to be adapted to deal with the fact that the data are not
simply measurements of each individual fish; full results are also given.
Using each method we calculate
the size-spectrum exponent $b$ separately for each year of data.
This gives us a time series of estimated annual values of $b$
with confidence intervals, one time series for each method. We then fit a
weighted linear regression (that uses the confidence intervals) to each
time series to see whether there is a
significant change ($p < 0.05$) in $b$ through time. For clarity we now
present the results which then motivate development of further methods.

\section{Results}

The linear regression results for each method (Figure~\ref{fig:nSeaFung/nSeaFungTrends}) show that
% From nSeaFungAnalysis.tex (which comes from ...Snw):
(a) the absolute estimates of
$b$ are highly dependent upon the method used, (b) the confidence intervals
can vary in size through time and are far broader for some methods than
for others,
% [Llin and MLE, but figure doesn't really imply that (hence I'd thought LCD
% before also), so leave details to Supp Mat.]
and (c) there is no significant change in $b$ through time when
using three of the methods
% [(namely LT, LTplus1 and MLE)],
yet a significant negative decline (steepening of the size spectrum)
when using the remaining five methods.
This latter point shows that
% five
% methods imply a steepening of the size spectrum over the thirty years of data,
% whereas three methods
% imply no change, demonstrating how
methodological differences can lead to
differing ecological conclusions.

% A weighted linear regression is fit to the resulting time series of estimates of $b$, to look for any trend (a weighted regression is used because we have estimates of the variance of each estimate of $b$, and these variances do vary between years).

% Have not defined these in main text - don't need to.
% Llin (Log-linear transform),
% LT (log transform),
% LTplus1 (log transform plus 1),
% LBmiz (logarithmic binning as done in the R package {\tt mizer} by
% \citealt{sba14}),

\subsection{Dealing with the bin structure of the data}

However, the aforementioned analyses implicitly ignore the fact that using the
minimum of a length
bin, and converting that to a single body mass, does not properly account for
the bin structure of the data (because the assigned length really represents any true
length within a bin). We demonstrate this issue using length-weight parameters
for relationship (\ref{lw}) for two example species included in \citet{blan05},
namely $\alpha_s= $0.0255 and $\beta_s=$2.7643 for Lemon Sole ({\it Microstomus
  kitt}) and $\alpha_s=$0.001 and $\beta_s=$3.4362 for Common Ling ({\it Molva
  molva}).
% Julia's file `{\tt weight length conversions final.R}' for Celtic Sea data

Figure \ref{fig:binFigs/binLW} shows how six 5-cm length bins get converted into
six body-mass bins via (\ref{lw}).
For clarity we use length bins of width 5~cm and use the midpoint of a
length bin to calculate the resulting body
mass (note that the IBTS data mostly uses 1-cm bins and
was analysed above using the minima of length bins rather than midpoints).
The resulting body-mass bins in Figure \ref{fig:binFigs/binLW} are not of equal
width because of the nonlinear nature of the length-weight relationships. And
the body-mass bins differ between species because of the species-specific
length-weight parameters. Such information regarding the body-mass bins was not
explicitly accounted for in the above analyses of the IBTS data; all methods
used a body mass obtained from applying (\ref{lw}) to the minimum of a length
bin.

Figure \ref{fig:binFigs/binLWlog} shows the consequence of using the resulting
body-mass values, converted from the midpoints of the length bins, when using
methods to fit size spectra that require further binning of the data. For
example, the LBbiom
(logarithmic binning to fit biomass size spectrum)
and LBNbiom (logarithmic binning with normalisation to fit biomass size
spectrum)
methods were used by \citet{jd05} and
\citet{pj05}, respectively, to calculate trends in size-spectra slopes for
earlier IBTS data. Explicitly shown is the assignment of the body-mass values
into bins with bin breaks of 4, 8, 16, 32, ..., as would be used in the LBbiom
and LBNbiom methods.

For example, consider blue length-bin number~5 that is highlighted in
Figure~\ref{fig:binFigs/binLW} for Lemon Sole and covers the range
30-35~cm. Say there are 20 individual Lemon Sole that have lengths in this
range. All 20 individuals would be represented by the length equal to the
midpoint of the length bin, namely 32.5~cm, which gets converted to a body-mass
value of 385~g (even though the 20 body masses really lie in the range
309-473~g). This value of 385~g lies within the $\logtwo$ bin that spans
256-512~g (Bin~5 in Figure~\ref{fig:binFigs/binLWlog}). Therefore, all 20
individuals get assigned a body mass of 384~g, which is the midpoint of the
$\logtwo$ bin. So, in the LBbiom and LBNbiom methods, all 20 individuals with
body masses in the range 309-473~g are assigned body masses of 384~g. The
uncertainty in the original lengths
(which is often an unavoidable consequence of measurement methods)
is not explicitly accounted for.
% Mike - think I'll leave text as is for now, especially if I can save words
% in the methods.
% If need to cut more words then see Mike's suggestions (17/7/18) on p10 and p12.

Furthermore, the length range 20-25~cm (Bin~3 in Figures~\ref{fig:binFigs/binLW}
and~\ref{fig:binFigs/binLWlog}) represents fish in the range 101-187~g.  The
midpoint of 22.5~cm gets converted to a body-mass value of 139~g, which lies
within the $\logtwo$ bin that spans 128-256~g
(Figure~\ref{fig:binFigs/binLWlog}). All individuals then get assigned a body
mass of 192~g (the midpoint of the $\logtwo$ bin), but this is \emph{greater}
than the body mass of any of the individuals (101-187~g).  A similar undesirable
situation occurs for Bin~6 (Figure~\ref{fig:binFigs/binLWlog}).  None of the
remaining bins yield counts in $\logtwo$ bins that accurately represent the
original length bins (since there is no reason that they should).

%Now, say there are 15~individual Lemon Sole in the length range 20-25~cm (Bin~3 in Figures~\ref{fig:binFigs/binLW} and~\ref{fig:binFigs/binLWlog}). The midpoint of this length bin, namely 22.5~cm, gets converted to a body-mass value of 139~g, though in reality the 15 body masses lie in the range 101-187~g. The value of 139~g lies within the $\logtwo$ bin that spans 128-256~g (Figure~\ref{fig:binFigs/binLWlog}). Therefore, all 15 individuals get assigned a body mass of 192~g, which is the midpoint of the $\logtwo$ bin. However, this body mass of 192~g is greater than the estimated body mass (101-187~g) of all 15 fish.

%Similarly, for Bin~6, the midpoint of the length range 35-40~cm gets converted to a body-mass value of 572~g, which lies within the $\logtwo$ bin that spans 512-1024~g with midpoint 768~g (Figure~\ref{fig:binFigs/binLWlog}). So all fish in Bin~6 get assigned a body mass of 768~g, but this is somewhat greater than the estimated body mass (473-684~g) of all individuals that Bin~6 represents. None of the remaining bins yield counts in $\logtwo$ bins that accurately represent the original length bins (since there is no reason that they should).

Similar situations also occur for the remaining four methods that involve
further binning of data.
%(Llin, LT, LTplus1 and LBmiz).
The LCD (logarithmic plotting of one minus the cumulative distribution) and
MLE (maximum likelihood estimate)
methods do not bin data and so do not
suffer from the extra binning described in Figure~\ref{fig:binFigs/binLWlog},
but do suffer from assigning the same body mass to all individuals in a
length bin.
It is not clear how the LCD method could be extended to account for the bin
structure. The
MLE method is the only method to accurately estimate $b$ and its confidence
intervals for unbinned data \citep{\MEE} --
this motivates our MLEbin
method, which adapts the MLE method to explicitly account for the uncertainty
inherent in binned data.

% END of from binLW.SNw




\subsection{Likelihood methods for binned data}

% Useful, if needed somewhere: MLEmid (which uses the midpoint of a body-mass bin as the body mass for all individual fish within that bin) and the MLEbin method (which explicitly accounts for the fact that the individuals within a bin can have a range of body sizes).

\subsubsection{MLEmid method}

The simplest way to deal with binned data in a likelihood framework is to just
use the midpoints of the body-mass bins, which we call the MLEmid (maximum
likelihood estimate based on mid-points) method. If there is, say, a count of 10
individuals assigned to the bin that spans the body-mass range 4-8~g, then all
10 individuals are assumed to have a body mass of 6~g (the midpoint, though
recall that for the IBTS data we followed \citealt{ffrr12} in using the minimum value for all eight methods,
including MLE). In
the Supplementary Material we derive the likelihood function for count data in
general and explain how this is used for the MLEmid method.

\subsubsection{MLEbin method}

The MLEbin (maximum likelihood estimate based on binning) method calculates the
MLE for $b$ by explicitly accounting for the fact that each count represent
values within its respective size bin. So for 10 individuals in the body-mass
bin 4-8~g, the MLEbin method uses a likelihood function that explicitly accounts
for the uncertainty in the 10 individual body masses. The body masses can be
anywhere in the range 4-8~g, rather than assuming that they are all 6~g as for
the MLEmid method. The likelihood function for the MLEbin method was derived as
equations (A.70) and (A.75) in \cite{\MEE}.

\subsubsection{Testing the MLEmid and MLEbin methods using simulated data}

We now simulate data to test the two methods. Following \citet{\MEE} we generate
10,000 simulated data sets, each containing 1,000 independent random numbers
(body masses)
drawn from the bounded power-law distribution (\ref{PLB}) with $b=-2$, $\xmin=1~$g
and $\xmax=1,000~$g. We bin each simulated data set to
represent the effect of collecting data in a binned form.
We then test how well
the MLEmid and MLEbin methods estimate $b=-2$ only using the information from
the binned data set (i.e. the bin definitions and the count in each bin), by
maximising the appropriate likelihood function and using the profile
likelihood-ratio test \citep{hm97} to calculate confidence intervals.

We test four binning types. Type `Linear 1' uses bins that all
have a constant bin width of 1~g and that span the data. The first bin starts at
the integer below the lowest data point.
Similarly, `Linear 5' and `Linear 10' use bins of constant widths 5~g
and 10~g, respectively. The binning type `2k' refers to bin widths that double in
size, span the data, and have bin breaks that are integer powers of 2 (i.e. are
at $2^k$~g where $k$ is an integer, giving bin breaks at 1, 2, 4, 8, 16, ...).

Figure~\fitMLEhists~shows that the MLEbin method clearly performs better than
the MLEmid method for all binning types (the histograms are centered around the
true value of $b=-2$); summary statistics are given in the Supplementary
Material. The MLEmid method always overestimates $b$, whereas the
MLEbin gives roughly the same results regardless of the binning (with a median
estimate of $b$ of -1.99 or -2.00).
% The MLEmid method performs best for the `Linear 1' binning type, but gets progressively worse as the constant bin width increases to 5 and then 10.

Somewhat surprisingly, the MLEmid method is biased even for the `Linear 1'
binning type. So, binning the data using bin breaks of 1, 2, 3, ..., 999, 1000~g
and then using the midpoint of a bin to represent all counts in that bin in a
likelihood context, gives a biased estimate of $b$ (99\% of the simulations
over-estimated $b$). This means that for, say, body-mass data that span the
range 1~g to 1~kg and are measured to a resolution of 1~g, ignoring the fact
that the binned data represent values in a range (i.e. using the MLEmid method)
will lead to biased values of the exponent $b$. The reason for this is the very
skewed nature of the power-law distribution. For example,
for the simulated data set
shown in Figure~\ref{fig:binFitting/exampleHists/histBinType1} of the
Supplementary Material, the
bin covering the range 1-2~g has a count of 528, but 348 (66\%) of these values
are actually $<1.5$~g; the MLEmid method assumes that all values equal
1.5~g. Whereas the MLEbin method, by definition, accounts for the fact that the
values within a bin are assumed to exhibit a power-law distribution.

% (rather than a uniform distribution as in the MLEmid method -- **think about to check this is true).

The MLEmid method performs particularly poorly with the `Linear 10' binning type
in Figure~\fitMLEhists(e), but, perhaps unexpectedly, the MLEbin method in
Figure~\fitMLEhists(f) leads to accurate estimation of $b$. So data collected at
a resolution of 10~g for which the MLEbin method is used to estimate $b$,
Figure~\fitMLEhists(f), give more accurate estimates of $b$ than collecting the
same data at a resolution of 1~g but using the MLEmid method to estimate $b$.
This shows that expending additional time and money to collect
high-resolution data
is not
worthwhile if the bin structure is unaccounted for -- it is better to collect
low-resolution data and use the MLEbin method to account for the bin structure.

The `2k' binning type performs slightly better than `Linear 10' when using the
MLEbin method (Figures~\fitMLEhists(f) and~(h)) because it has higher resolution
for the data-rich lower bins (bin breaks are 1, 2, 4, 8, ..., compared to 1, 11,
21, 31, ..., for `Linear 10').

The results regarding confidence intervals (Figure~\fitMLEconfs) give a similar
conclusion. For the MLEmid method and `Linear 1' binning type, only 46\% of the
95\% confidence intervals contain the true value of $b$, less than half of the
desired observed coverage of 95\%. The results are worse as the bin width
increases (`Linear 5' and `Linear 10') and for bin widths that double
(`2k'). However, the MLEbin gives reliable confidence intervals (observed
coverage of 94\% or 95\%) regardless of binning type.
%The intervals get slightly
%wider as the bin width increases (from Figure~\fitMLEconfs(b) to (d) to (f)) as
%the resolution of the data gets poorer. **Check that with the simulation results
%- maybe it's not exactly major.

Thus, these results demonstrate that for binned data the MLEmid method is
inaccurate and the MLEbin method should be used, even for seemingly narrow bins.

\subsubsection{The MLEbins method  --  accounting for species-specific body-mass bins}

The IBTS data has species-specific length-weight
relationships, such that the
body-mass bins are different for each species, as occurs for many data sets;
e.g.~\citet{mnwb18} had lengths of 1-cm resolution that were
converted into body masses (and then used for the LBNbiom method).
To account for this we extend the MLEbin
method to the MLEbins method, using a multinomial
log-likelihood function \citep{lawl03}; see Supplementary Material.

\subsection{MLEbins method applied to IBTS data}

To apply the MLEbins method to the IBTS data we first need to calculate
species-specific body-mass bins that result from the species-specific
length-weight coefficients (as for Figure~\ref{fig:binFigs/binLW}).
%Adapted From nSeaFungMLEbin.Snw/.tex:
Figure~\ref{fig:nSeaFung/nSeaFungBins1} shows how the 1-cm (or 0.5-cm) length
bins for 45 of the species get converted to body-mass bins. The conversions are
different for each species because of the different values of the length-weight
coefficients. Thus, the vertical bars are analogous to those in
Figure~\ref{fig:binFigs/binLW}. Figure~\ref{fig:nSeaFung/nSeaFungBins1} shows
that even though the length bins are only 1~cm wide, the resulting individual
body-mass bins can be quite large; the largest across all species is 832~g for
Atlantic Cod (\emph{Gadus morhua}) [see Supplementary Material].

In Figure~\ref{fig:ISD-1999.main} we introduce a plotting technique
that visualises the binning structure of the data and the fit of the ISD
(from using the MLEbins method).
The fitted PLB model shows
a reasonable fit to the 1999 data (the red curve passing through the grey
rectangles) for body masses $<100$~g; for larger body masses the PLB model
generally over-estimates the numbers of individuals. This pattern is seen for
many of the years (see Supplementary Material for figures)
and may potentially be a consequence of fishing pressure.


Using the species-specific MLEbins method applied to each year of the IBTS data, we obtain
consistently higher estimates of $b$ than for the original MLE method
(Figure~\ref{fig:nSeaFung/nSeaFungCompareTrends}).
The differing estimates of $b$ show that the likelihood method did need to be
properly adapted (the MLEbins method) to deal with the binned nature of the
data.
In this case, both methods conclude that there is no significant temporal trend
in $b$ through time (see Supplementary Material). Nevertheless, the results show
the potential for different conclusions about community structure to be reached
if the bin structure of data is not properly considered.

\subsection{Recommendations}

In Table~\ref{tab:recommend} we recommend fitting methods based on the type
and resolution of data available.
If data are available in unbinned form (e.g.~\citealt{tdss15})
then the MLE method can be used, but this still depends on the accuracy of the raw
data. Given that unbinned data are
essentially (to some resolution) binned data,
given the results in Figures~\fitMLEhists\ and~\fitMLEconfs\ we recommend
testing the MLE versus MLEbin method for the particular bin structure and range
of a data set.
Or one could just use the MLEbin method, with the bin structure
defined by the accuracy of the measurements.
Our methods may need to be tailored for explicit data sets depending on the
survey design -- in Table~\ref{tab:recommend} we give an example concerning
monitoring of Pacific reefs.

\section{Discussion}

We have analysed a rich 30-year data set from the North Sea, and found that
ecological conclusions regarding trends in size-spectra exponent can depend on
the method used to estimate the exponent. Furthermore, the likelihood method
needs to properly account for the binned nature of data (the MLEbin method,
adapted as the MLEbins method for the IBTS data). Using simulated data, even for
seemingly high-resolution data the likelihood method that does not account for
the binning (the MLEmid method) is inaccurate, whereas MLEbin is accurate.
For data, such as the IBTS data, where lengths have to be converted to body
masses using species-specific length-weight coefficients, we have derived
the MLEbins method to estimate $b$ and developed a technique for plotting the
data and the resulting fit.

Following \citet{\MEE} we have used a bounded power-law distribution for the ISD
because power laws are commonly used models for size spectra \citep{pd78, bd92}
and alternative models of size-structured predator-prey dynamics predict
that the aggregate community ISD may indeed be
close to a power law \citep{ab06}.
% , the distributions of individual species may
% not be \citep{hab11, jga14, lpk14}. For the species-specific MLEbins method we
%have necessarily assumed that ISD's for species are power laws with a common
% exponent $b$ (specifically equation (\ref{Pindiv}) in the Supplementary
% Material).
Future work could focus on relaxing the
assumption of a power-law distribution
if other distributions are deemed feasible \citep{\MEE}.
Using a bounded (rather than unbounded) distribution avoids extrapolation
beyond the range of the data -- this regrettably prevents estimation
of the population density of Loch Ness monsters (as per \citealt{sk72}).
Since power-law distributions have many applications in ecology \citep{weg08},
our simulation results suggest that bin structure may need to be accounted for
in other contexts, and our novel plotting method
(Figure~\ref{fig:ISD-1999.main}) would be applicable.

Following \citet{pj05} and \citet{blan05} we used a (minimum) cut-off of 4~g to
analyse the IBTS data. This was necessarily done by converting the minima of the
length bins into species-specific body-mass bins, and then only using the
resulting bins with minima $\geq$4~g.  For a resulting body-mass bin of
3-5~g, for example, this approach removes some individuals $\geq$4~g, even though we wish to
impose a minimum of 4~g. This may be unavoidable, although note that the
likelihood function for the MLEbins method properly assumes no data (rather than no
individuals) $<$5~g for such a species with lowest bin 3-5~g [$w_{s,1} = 5$~g
for this species in (\ref{loglikbin}) in the Supplementary Material].
Also note that we have estimated $\xmax$ separately for each year although it
could be set to a fixed maximum value across all years.

We have used the IBTS data to motivate development of widely applicable
methods. We caution against using our current results for management purposes
concerning the North Sea, because we have not fully explored assumptions (such
as those just discussed) and have used all the sampled species rather than only
a well-sampled subset (as recommended by \citealt{jd05}). Also, the power
analyses conducted by \citet{jd05} on earlier IBTS data could be repeated, in
particular to determine the lower cut-off value that would provide the greatest
power to detect trends in the size-spectrum exponent. Note that \citet{jd05}
were restricted to such cut-offs being powers of two because of the use of the
LBbiom method, but this would not be necessary with the MLEbins method.


Note that the confidence intervals calculated for the IBTS data
from the MLEbins method
are not strictly true confidence intervals, because the commonly-used unit of `numbers
of individuals caught per hour of trawling' is arbitrary. If numbers were
instead expressed as `per day of trawling' then all counts would be divided by
24, and
the resulting confidence intervals would be larger than those in
Figure~\ref{fig:nSeaFung/nSeaFungCompareTrends}.
True confidence intervals could be obtained by using the absolute numbers of
individual fish caught (rather than numbers per hour of
trawling) summed across all areas.
This would change the width of the confidence intervals, but not the MLE for
$b$, which is unaffected by any change in units.

All results, figures and tables presented here are reproducible from the
supplied R code. The code is documented and functionalised with a view to it
being used by others to fit size spectra to their data, using our recommended
and tested likelihood methods that properly account for the bin structure of the data.

% Refs in Supp but not in main text: ba02, edwa07, efbj12 - now done manually

%TC:ignore
\section*{Supplementary Material}

The following supplementary material is included:
(A)~document containing further methodological details, including derivations
of appropriate likelihood functions and further results for the IBTS and
the simulated data;
(B)~code repository such that all results can be reproduced and users
can apply the methods to their own data.
AME is creating an R package from the code for even easier use by others.

\section*{Acknowledgements}

We thank Carrie Holt, Brooke Davis, Jaclyn Cleary and Jennifer Boldt for useful
discussions.
%TC:endignore


%TC:incbib

\begin{thebibliography}{}

\bibitem[\'Alvarez et~al.(2016)]{amln16}
\'Alvarez, E., Mor\'an, X. A.~G., L\'opez-Urrutia, A., and Nogueira, E. (2016).
\newblock {Size-dependent photoacclimation of the phytoplankton community in
  temperate shelf waters (southern Bay of Biscay)}.
\newblock {\em Mar.~Ecol.~Prog.~Ser.}, {\bf 543}, 73--87.

\bibitem[Andersen \& Beyer(2006)]{ab06}
Andersen, K.~H. \& Beyer, J.~E. (2006).
\newblock Asymptotic size determines species abundance in the marine size
  spectrum.
\newblock {\em Am.~Nat.}, {\bf 168}, 54--61.

\bibitem[Blanchard et~al.(2005)]{blan05}
Blanchard, J.~L., Dulvy, N.~K., Jennings, S., Ellis, J.~R., Pinnegar, J.~K.,
  Tidd, A., \& Kell, L.~T. (2005).
\newblock Do climate and fishing influence size-based indicators of {C}eltic
  {S}ea fish community structure?
\newblock {\em ICES J.~Mar.~Sci.}, {\bf 62}, 405--411.

\bibitem[Boldt et~al.(2012)]{bblhw12}
Boldt, J.~L., Bartkiw, S.~C., Livingston, P.~A., Hoff, G.~R., \& Walters,
  G.~E. (2012).
\newblock Investigation of fishing and climate effects on the community size
  spectra of eastern {B}ering {S}ea fish.
\newblock {\em Trans.~Am.~Fish.~Soc.}, {\bf 141}, 327--342.

\bibitem[Boudreau \& Dickie(1992)]{bd92}
Boudreau, P.~R. \& Dickie, L.~M. (1992).
\newblock Biomass spectra of aquatic ecosystems in relation to fisheries yield.
\newblock {\em Can. J. Fish. Aquat. Sci.}, {\bf 49}, 1528--1538.

\bibitem[Daan et~al.(2005)]{dgpr05}
Daan, N., Gislason, H., Pope, J.~G., \& Rice, J.~C. (2005).
\newblock Changes in the {N}orth {S}ea fish community: evidence of indirect
  effects of fishing?
\newblock {\em ICES J.~Mar.~Sci.}, {\bf 62}, 177--188.

\bibitem[Dulvy et~al.(2004)]{dpmg04}
Dulvy, N.~K., Polunin, N. V.~C., Mill, A.~C., \& Graham, N. A.~J. (2004).
\newblock Size structural change in lightly exploited coral reef fish
  communities: evidence for weak indirect effects.
\newblock {\em Can. J. Fish. Aquat. Sci.}, {\bf 61}, 466--475.

\bibitem[Edwards(2011)]{edwa11}
Edwards, A.~M. (2011).
\newblock Overturning conclusions of {{L}}\'evy flight movement patterns by
  fishing boats and foraging animals.
\newblock {\em Ecology}, {\bf 92}, 1247--1257.

\bibitem[Edwards et~al.(2017)]{erpbb17}
Edwards, A.~M., Robinson, J. P.~W., Plank, M.~J., Baum, J.~K., \& Blanchard,
  J.~L. (2017).
\newblock Testing and recommending methods for fitting size spectra to data.
\newblock {\em Meth.~Ecol.~Evol.}, {\bf 8}, 57--67.

\bibitem[{European~Commission}(2010)]{europe10}
{European~Commission} (2010).
\newblock {Commission decision of 1 September 2010 on criteria and
  methodological standards on good environmental status of marine waters}.
\newblock {\em Off. J. Eur. Union}, {\bf L232}, 14--24.

\bibitem[Fung et~al.(2012)]{ffrr12}
Fung, T., Farnsworth, K.~D., Reid, D.~G., \& Rossberg, A.~G. (2012).
\newblock Recent data suggest no further recovery in {N}orth {S}ea {L}arge
  {F}ish {I}ndicator.
\newblock {\em ICES J.~Mar.~Sci.}, {\bf 69}, 235--239.

\bibitem[Heenan et~al.(2016)]{heen16}
Heenan, A., Gorospe, K., Williams, I., Levine, A., Maurin, P., Nadon, M.,
  Oliver, T., Rooney, J., Timmers, M., Wongbusarakum, S., \& Brainard, R.
  (2016).
\newblock Ecosystem monitoring for ecosystem-based management: using a
  polycentric approach to balance information trade-offs.
\newblock {\em J.~Appl.~Ecol.}, {\bf 53}, 699--704.

\bibitem[Hilborn \& Mangel(1997)]{hm97}
Hilborn, R. \& Mangel, M. (1997).
\newblock {\em The Ecological Detective: Confronting Models with Data}.
\newblock Vol.~28, Monographs in Population Biology, Princeton University
  Press, New Jersey.

\bibitem[ICES(2015)]{ices15}
ICES (2015).
\newblock {Manual for the International Bottom Trawl Surveys}.
\newblock Series of ICES Survey Protocols SISP 10 -- IBTS IX. 86~p.

\bibitem[Jennings(2005)]{jenn05}
Jennings, S. (2005).
\newblock Indicators to support an ecosystem approach to fisheries.
\newblock {\em Fish Fish.}, {\bf 6}, 212--232.

\bibitem[Jennings \& Dulvy(2005)]{jd05}
Jennings, S. \& Dulvy, N.~K. (2005).
\newblock Reference points and reference directions for size-based indicators
  of community structure.
\newblock {\em ICES J.~Mar.~Sci.}, {\bf 62}, 397--404.

\bibitem[Lawless(2003)]{lawl03}
Lawless, J.~F. (2003).
\newblock {\em Statistical Models and Methods for Lifetime Data}.
\newblock 2nd ed., Wiley series in Probability and Statistics, Wiley, New
  Jersey.

\bibitem[Maxwell \& Jennings(2006)]{mj06}
Maxwell, T. A.~D. \& Jennings, S. (2006).
\newblock Predicting abundance-body size relationships in functional and
  taxonomic subsets of food webs.
\newblock {\em Oecologia}, {\bf 150}, 282--290.

\bibitem[McCoy et~al.(2016)]{mcco16}
McCoy, K., Heenan, A., Asher, J., Ayotte, P., Gorospe, K., Gray, A., Lino, K.,
  Zamzow, J., \& Williams, I. (2016).
\newblock {Pacific Reef Assessment and Monitoring Program data report.
  Ecological monitoring 2015 -- reef fishes and benthic habitats of the main
  Hawaiian Islands, Northwest Hawaiian Islands, Pacific Remote Island Areas,
  and American Samoa}.
\newblock PIFSC Data Report DR-16-002.

\bibitem[Mindel et~al.(2018)]{mnwb18}
Mindel, B.~L., Neat, F.~C., Webb, T.~J., \& Blanchard, J.~L. (2018).
\newblock Size-based indicators show depth-dependent change over time in the
  deep sea.
\newblock {\em ICES J.~Mar.~Sci.}, {\bf 75}, 113--121.

\bibitem[Monnahan et~al.(2016)]{monn16}
Monnahan, C.~C., Ono, K., Anderson, S.~C., Rudd, M.~B., Hicks, A.~C.,
  Hurtado-Ferro, F., Johnson, K.~F., Kuriyama, P.~T., Licandeo, R.~R., Stawitz,
  C.~C., Taylor, I.~G., \& Valero, J.~L. (2016).
\newblock The effect of length bin width on growth estimation in integrated
  age-structured stock assessments.
\newblock {\em Fish.~Res.}, {\bf 180}, 103--112.

\bibitem[Piet \& Jennings(2005)]{pj05}
Piet, G.~J. \& Jennings, S. (2005).
\newblock Response of potential fish community indicators to fishing.
\newblock {\em ICES J.~Mar.~Sci.}, {\bf 62}, 214--225.

\bibitem[Platt \& Denman(1978)]{pd78}
Platt, T. \& Denman, K. (1978).
\newblock The structure of pelagic marine ecosystems.
\newblock {\em Rapp. P.-v. R{{\'e}}un. Cons. Int. Explor. Mer}, {\bf 173},
  60--65.

\bibitem[Rice \& Gislason(1996)]{rg96}
Rice, J. \& Gislason, H. (1996).
\newblock Patterns of change in the size spectra of numbers and diversity of
  the {N}orth {S}ea fish assemblage, as reflected in surveys and models.
\newblock {\em ICES J.~Mar.~Sci.}, {\bf 53}, 1214--1225.

\bibitem[Rice et~al.(2012)]{rice12}
Rice, J., Arvanitidis, C., Borja, A., Frid, C., Hiddink, J.~G., Krause, J.,
  Lorance, P., Ragnarsson, S.~A., Sk\"old, M., Trabucco, B., Enserink, L., \&
  Norkko, A. (2012).
\newblock {Indicators for sea-floor integrity under the European Marine
  Strategy Framework Directive}.
\newblock {\em Ecol.~Indic.}, {\bf 12}, 174--184.

\bibitem[Richards et~al.(2011)]{rwnz11}
Richards, B.~L., Williams, I.~D., Nadon, M.~O., \& Zgliczynski, B.~J. (2011).
\newblock A towed-diver survey method for mesoscale fishery-independent
  assessment of large-bodied reef fishes.
\newblock {\em Bull.~Mar.~Sci.}, {\bf 87}, 55--74.

\bibitem[Scott et~al.(2014)]{sba14}
Scott, F., Blanchard, J.~L., \& Andersen, K.~H. (2014).
\newblock \emph{mizer}: an {R} package for multispecies, trait-based and
  community size spectrum ecological modelling.
\newblock {\em Meth.~Ecol.~Evol.}, {\bf 5}, 1121--1125.

\bibitem[Sheldon \& Kerr(1972)]{sk72}
Sheldon, R.~W. \& Kerr, S.~R. (1972).
\newblock {The population density of monsters in Loch Ness}.
\newblock {\em Limnol.~Oceanogr.}, {\bf 17}, 796--798.

\bibitem[Shin et~al.(2005)]{srjfg05}
Shin, Y.-J., Rochet, M.-J., Jennings, S., Field, J.~G., \& Gislason, H.
  (2005).
\newblock Using size-based indicators to evaluate the ecosystem effects of
  fishing.
\newblock {\em ICES J.~Mar.~Sci.}, {\bf 62}, 384--396.

\bibitem[Stasko et~al.(2016)]{ssmarp16}
Stasko, A.~D., Swanson, H., Majewski, A., Atchison, S., Reist, J., \& Power,
  M. (2016).
\newblock {Influences of depth and pelagic subsidies on the size-based trophic
  structure of Beaufort Sea fish communities}.
\newblock {\em Mar.~Ecol.~Prog.~Ser.}, {\bf 549}, 153--166.

\bibitem[Taniguchi et~al.(2014)]{tfp14}
Taniguchi, D. A.~A., Franks, P. J.~S., \& Poulin, F.~J. (2014).
\newblock Planktonic biomass size spectra: an emergent property of
  size-dependent physiological rates, food web dynamics, and nutrient regimes.
\newblock {\em Mar.~Ecol.~Prog.~Ser.}, {\bf 514}, 13--33.

\bibitem[Thorpe et~al.(2015)]{tllcj15}
Thorpe, R.~B., {Le~Quesne}, W. J.~F., Luxford, F., Collie, J.~S., \& Jennings,
  S. (2015).
\newblock Evaluation and management implications of uncertainty in a
  multispecies size-structured model of population and community responses to
  fishing.
\newblock {\em Meth.~Ecol.~Evol.}, {\bf 6}, 49--58.

\bibitem[Trebilco et~al.(2015)]{tdss15}
Trebilco, R., Dulvy, N.~K., Stewart, H., \& Salomon, A.~K. (2015).
\newblock The role of habitat complexity in shaping the size structure of a
  temperate reef fish community.
\newblock {\em Mar.~Ecol.~Prog.~Ser.}, {\bf 532}, 197--211.

\bibitem[{van~Gemert} \& Andersen(2018)]{va18}
{van~Gemert}, R. \& Andersen, K.~H. ({2018}).
\newblock Implications of late-in-life density-dependent growth for fishery
  size-at-entry leading to maximum sustainable yield.
\newblock {\em ICES J.~Mar.~Sci.}, {\bf 75}, 1296--1305.

\bibitem[White et~al.(2007)]{weke07}
White, E.~P., Ernest, S. K.~M., Kerkhoff, A.~J., \& Enquist, B.~J. (2007).
\newblock Relationships between body size and abundance in ecology.
\newblock {\em Trends Ecol.~Evol.}, {\bf 22}, 323--330.

\bibitem[White et~al.(2008)]{weg08}
White, E.~P., Enquist, B.~J., \& Green, J.~L. (2008).
\newblock On estimating the exponent of power-law frequency distributions.
\newblock {\em Ecology}, {\bf 89}, 905--912.

\bibitem[Woodson et~al.(2018)]{wsj18}
Woodson, C.~B., Schramski, J.~R., \& Joye, S.~B. (2018).
\newblock A unifying theory for top-heavy ecosystem structure in the ocean.
\newblock {\em Nat. Commun.}, {\bf 9}, 23.

\bibitem[Zhang et~al.(2018)]{zcxxr18}
Zhang, C., Chen, Y., Xu, B., Xue, Y., \& Ren, Y. (2018).
\newblock Evaluating fishing effects on the stability of fish communities using
  a size-spectrum model.
\newblock {\em Fish.~Res.}, {\bf 197}, 123--130.

\bibitem[Zhou et~al.(in press)]{zhou19}
Zhou, S., Kolding, J., Garcia, S.~M., Plank, M.~J., Bundy, A., Charles, A.,
  Hansen, C., Heino, M., Howell, D., Jacobsen, N.~S., Reid, D.~G., Rice, J.~C.,
  \& {van~Zwieten}, P. A.~M. ({\protect{in press}}).
\newblock Balanced harvest: concept, policies, evidence, and management
  implications.
\newblock {\em Rev.~Fish Biol.~Fish.},
  https://doi.org/10.1007/s11160-019-09568-w.

\end{thebibliography}

% \bibliography{../../../latex/abbrev5}

\clearpage

% From nSeaFungAnalysis.Snw. If replace here then copy first as need to
%  change some things manually.
% Ordering is by year then specCode, which won't be obvious once have
%  species names.
% 105814 has two names. Go with Scyliorhinus canicula (Smallspotted Catshark
%  [Sharks of the world, Compagno et al.])
% 154675: Lumpenus lampretaeformis  (Snakeblenny)
% 274304: Microchirus variegatus (Thickback Sole)

% Species names
% latex table generated in R 3.4.3 by xtable 1.8-2 package
% Tue May 15 15:12:05 2018
\begin{table}[tp]
\centering
\caption{First six and last six rows of the IBTS data to illustrate the
  information available. Each row represents a unique combination of year,
  species and length class. `Number' gives the number of individuals per hour of
  trawling observed for that combination, and can be non-integer because counts
  of individual fish are scaled by tow duration. Fish lengths are assigned into
  length classes, where the value of `Length class' (cm) is the minimum value of
  the 1-cm length bin (or 0.5-cm length bins for Atlantic Herring and European
  Sprat). Parameters $\alpha_s$ and $\beta_s$ are the length-weight coefficients
  for species $s$ from \citet{ffrr12}, and `Body mass' (g) is the resulting
  estimated body mass for an individual of that species and length
  class. `Biomass' (g~h$^{-1}$) is the total biomass caught per hour of trawling
  for each row. Example species are Smallspotted Catshark ({\it Scyliorhinus
    canicula}), Snakeblenny ({\it Lumpenus lampretaeformis}) and Thickback Sole
  ({\it Microchirus variegatus}). The full data set has 42,298 rows.}
\label{tab:egData}
\begin{tabular}{rrrrrrrr}
  \hline
  Year & Species & Length & Number & $\alpha_s$ & $\beta_s$ & Body & Biomass\\
       & & class (cm) & (h$^{-1}$) & & & mass (g) & (g~h$^{-1}$) \\
  \hline
  1986 & Smallspotted Catshark & 45 & 0.007 & 0.0031 & 3.0290 & 315.46 & 2.25 \\
  1986 & Smallspotted Catshark & 46 & 0.007 & 0.0031 & 3.0290 & 337.17 & 2.41 \\
  1986 & Smallspotted Catshark & 50 & 0.007 & 0.0031 & 3.0290 & 434.05 & 3.10 \\
  1986 & Smallspotted Catshark & 52 & 0.029 & 0.0031 & 3.0290 & 488.81 & 14.33 \\
  1986 & Smallspotted Catshark & 53 & 0.011 & 0.0031 & 3.0290 & 517.84 & 5.65 \\
  1986 & Smallspotted Catshark & 54 & 0.011 & 0.0031 & 3.0290 & 548.00 & 6.18 \\
  $\cdot \cdot \cdot$ & $\cdot \cdot \cdot$ & $\cdot
  \cdot \cdot$ & $\cdot \cdot \cdot$ & $\cdot \cdot \cdot$ & $\cdot \cdot \cdot$
               & $\cdot \cdot \cdot$ & $\cdot \cdot \cdot$ \\
  2015 & Snakeblenny & 34 & 0.028 & 0.0244 & 2.0439 & 32.93 & 0.92 \\
  2015 & Thickback Sole & 8 & 0.013 & 0.0080 & 3.1410 & 5.49 & 0.07 \\
  2015 & Thickback Sole & 14 & 0.039 & 0.0080 & 3.1410 & 31.85 & 1.24 \\
  2015 & Thickback Sole & 15 & 0.052 & 0.0080 & 3.1410 & 39.55 & 2.05 \\
  2015 & Thickback Sole & 16 & 0.065 & 0.0080 & 3.1410 & 48.44 & 3.15 \\
  2015 & Thickback Sole & 17 & 0.013 & 0.0080 & 3.1410 & 58.60 & 0.76 \\
  \hline
\end{tabular}
\end{table}

\begin{table}[tp]
\centering
\caption{Recommendation table for how to calculate the size-spectrum exponent,
  $b$, of the individual size distribution for body masses (equation \ref{PLB}),
  depending on the type and resolution of the data.}
\label{tab:recommend}
\small{
\begin{tabular}{L{3 cm}L{3 cm}L{5 cm}L{4 cm}}
\hline
Data type & Resolution & Issue & Solution \\
\hline
  Body mass & Individual body masses & Accuracy of measurements may require them
                                       to be considered as bins
          & Simulate data as in Figure~\fitMLEhists~to see if MLE method
            sufficient or MLEbin is needed\\
Body mass & Individuals assigned to common body-mass bins & Data are counts in bins
                               & Use MLEbin method\\
Length & Individual lengths & Lengths have to be converted to body masses
       & Calculate individual body masses and proceed as per first row\\
Length & Individuals assigned to length bins &
         Resulting converted body-mass bins are species specific &
         Use MLEbins method\\
Lengths from underwater visual census (e.g.~point counts or belt transects)
       & Individual lengths
       & May need to standardise counts when combining different survey designs
         (e.g.~\citealt{heen16}); surveys may cover different sized areas or
         count small and large organisms differently
       & Need to adapt MLEbins likelihood function to account for particular
         survey design \\
       %& Individual lengths of small fish within 5~m of transect,
       %      medium fish within **~m of transect, large fish within 30~m
       %      of transect & Need to standardise counts & **See section A.5,
       %       need to flesh out likelihood function\\
% Resolution: 'Individual lengths, where small and large organisms are counted in di%fferent survey areas'
%Issue: 'Counts need to be corrected by survey area'
\hline
\end{tabular}
}
\end{table}

\renewcommand{\baselinestretch}{1.1}

% was dataCode/northSeaFung/nSeaFungTrends
\onefigshift{nSeaFung/nSeaFungTrends}{For the IBTS data, each of eight methods is
  used to estimate the exponent $b$ (circles) and associated 95\% confidence
  interval (vertical bars) for each year. The fit of a weighted linear
  regression with 95\% confidence interval is shown in red if the trend
  can be considered statistically significant from 0 ($p<0.05$) and in grey if
  it cannot be ($p \geq 0.05$).
  The y-axes are the same for (d)-(h). For (g) and (h) the confidence
  intervals are too narrow to be seen. See Supplementary Material for full
  details of methods.}
  % Statistical properties of trends are
  % given in Table~\ref{tab:trendRes}.}

% was dataCode/binLW

\onefigshift{binFigs/binLW}{Example length-weight relationships (solid curves)
  for Lemon Sole ({\it Microstomus kitt}) in blue and Common Ling ({\it Molva
    molva}) in red. Horizontal bars show six example 5-cm length bins, the same
  bins for both species. Using the length-weight relationships, these length
  bins convert to the body-mass bins shown by the vertical bars, giving six bins
  of unequal width for each species. The conversion for the 30-35~cm length bin
  is explicitly shown as dotted lines for the endpoints and solid lines for the
  midpoint of 32.5~cm. For Lemon Sole, the 30-35~cm bin with midpoint 32.5~cm
  converts to a body-mass bin of 309-473~g (values
  rounded to nearest~g for clarity), and for Common Ling the converted body-mass
  bin is 119-202~g.
  Note that converting the midpoint of 32.5~cm gives respective body masses of
  385~g and 157~g, which are slightly
  less than the respective midpoints of the new body-mass bins (391~g and
  161~g) because of the nonlinearity of the length-weight relationships.}
% sprintf("%.1f", round(massBinsStart[1,5], dig=1))

\clearpage
% was dataCode/binLWlog

\onefig{binFigs/binLWlog}{Demonstration of how binned body-mass values
  (blue bars) get
  assigned to logarithmic size-class bins (black and grey bars),
  as occurs for the LBbiom and LBNbiom
  methods that prescribe the logarithmic bins. Blue vertical bars show the
  resulting six converted body-mass bins for Lemon Sole from
  Figure~\ref{fig:binFigs/binLW}. The black and grey vertical bars show
  body-mass bins that have equal width on a logarithmic scale. Explicitly these
  are at $4, 8, 16, 32, ...$, as used in the LBbiom and LBNbiom methods. The
  short horizontal purple lines show the midpoints of those bins. Each panel,
  labelled Bin~1 to Bin~6, highlights one blue body-mass bin. For each blue bin,
  the blue arrow starts from the converted value of the midpoint of the original
  length bin (see Figure~\ref{fig:binFigs/binLW} for Bin~5), and ends at the
  purple midpoint of the $\logtwo$ bin in which the converted value lies. The
  blue dashed lines indicate that all body masses within the blue bin get
  assigned to the same purple midpoint value. This ignores uncertainties in
  body masses, and for Bins~3 and~6 assigns completely incorrect values
  (greater than the body mass of any individual within the bin). But
  such values get
  used in the LBbiom or LBNbiom methods.}

\clearpage

\thispagestyle{empty}

% \fitMLEhists
\onefig{binFitting/fitMLEhists}{Histograms of estimated exponent $b$ for 10,000
  simulated data sets, each of which contains 1,000 independent random numbers
  drawn from a bounded power-law distribution with $b=-2$, $\xmin=1$ and
  $\xmax=1,000$.
%Each simulated data set is binned in four ways, with linear bins of width 1 (panels a and b), 5 (c, d), and 10 (e, f), or bin widths that double in size (g, h).
Panels (a)-(f) use linear bins of constant bin width 1, 5 or 10 (binning type
`Linear 1' etc.), and panels (g)-(h) use bin widths that double in size (binning
type `2k'). Each binned data set is fitted using either the MLEmid method (a, c,
e, g) or the MLEbin method (b, d, f, h). Vertical red lines indicate the known
value of $b=-2$.}

\clearpage
\thispagestyle{empty}

% \fitMLEconfs
\onefigshift{binFitting/fitMLEconfs}{Confidence intervals (horizontal lines) of
  $b$ obtained for each combination of data-binning and MLE method for
  subsamples of the 10,000 simulated data sets (with $\xmax=1,000$) used in
  Figure~\fitMLEhists. Panels are the same as in Figure~\fitMLEhists. For each
  numbered subsample on the y-axis, the 95\% confidence interval of $b$ obtained
  using the respective method is plotted as a horizontal line, which is coloured
  grey if the interval includes the true value of $b=-2$ (given by the vertical
  red line) or blue if it does not. Simulations are sorted in ascending order of
  their lower bound. The percentage for each method gives the observed coverage,
  namely the percentage of all 10,000 simulated data sets for which the 95\%
  confidence interval contains the true value of $b$; by definition, this should
  ideally be 95\%. Horizontal axes are the same for all panels here and in
  Figure~\fitMLEhists.}

%OLD part 2:
%(c) Confidence intervals (horizontal lines) of $b$ obtained using the MLEmid  method for subsamples of the 10,000 simulated data sets. For each numbered subsample on the y-axis, the 95\% confidence interval of $b$ is plotted as a horizontal line, which is coloured grey if the interval includes the true value of $b=-2$ (given by the vertical red line) or blue if it does not. The simulations are sorted in ascending order of the lower bound. The percentage shown is the observed coverage, namely the percentage of all 10,000 simulated data sets for which the 95\% confidence interval contains the true value of $b$; by definition, this should ideally be 95\%. (d) As for (c) but using the MLEbin method.}

\clearpage

\onefig{nSeaFung/nSeaFungBins1}{Vertical bars show the body-mass bins that
  result from converting the 1-cm length bins into body-mass bins using the
  species-specific length-weight parameters. These are the first 45 species of
  the IBTS data, ordered by the maximum value of the highest bin. Atlantic
  Herring and European Sprat use 0.5-cm length bins (indicated by~$\times$). Any
  vertical gaps between vertical bars indicate that there were no body masses of
  that species in that range; i.e.~there were no individuals in the length
  bin(s) that would result in the body-mass bin(s) that would fill the gap. For
  example, there were no $\sim$30~g sprat (code 126425). Two example species are
  shown in red. They have similar ranges of body masses, but Moustache Sculpin
  (code 127205, \emph{Triglops murrayi}) is short and fat with lengths from
  8-15~cm and length-weight relationship $w = 0.0088~l^{3.00}$, whereas
  Snakeblenny (code 154675, \emph{Lumpenus lampretaeformis}) is long and skinny
  with lengths from 13-37~cm and $w = 0.0244~l^{2.04}$, and consequently the
  1-cm length bins translate to much narrower body-mass bins than for Moustache
  Sculpin. Equivalent figures for the remaining 90 species are given in the
  Supplementary Material.}


\clearpage
\thispagestyle{empty}
\begin{figure}[tp]
\begin{center}
\vspace{-10mm}
\epsfbox{../../fitting-size-spectra/codeBinStructure/nSeaFung/IBTS-ISDfigs/IBTS-ISD1999.eps}
\end{center}
\vspace{-5mm}
% Change caption in Supp Infor for 1986 also, if edit here; not identical though
\caption{Suggested plots of individual size distribution and MLEbins fit
            (red solid curve)
            with 95\% confidence intervals (red dashed curves, may be hard to
            see) for IBTS data in 1999. The y-axis is linear in (a) and logarithmic
            in (b).
            For each bin, the horizontal green line shows
            the range of body sizes,
            with its value on the y-axis corresponding to the
            total number of individuals (per hour of trawling) in bins whose
            minima are $\geq$ the
            bin's minimum.
            The vertical span of each grey rectangle shows the
            possible range of the number of individuals with body mass
            $\geq$ the body mass of individuals in that bin (horizontal span
            is the same as for the green lines).
            The text in (a) gives the year,
            the MLE for the size-spectrum exponent $b$, and the sample size $n$.
            See Supplementary Material for
            full details and results for all years.}
\label{fig:ISD-1999.main}
\end{figure}
\clearpage

\begin{figure}[tp]
\begin{center}
\epsfxsize=6in
\epsfbox{nSeaFungCompareTrendsCol.eps}
  \end{center}
  \caption{Comparison of the annual estimates of
  $b$ for the IBTS data using the original MLE method and the MLEbins method,
  showing medians and 95\% confidence intervals (each year is slightly offset
  to clearly show the confidence intervals).}
\label{fig:nSeaFung/nSeaFungCompareTrends}
\end{figure}

\clearpage

%TC:ignore

% \section*{Supplementary Material}

% \vspace{-10mm}
\part{Supplementary Material -- Extended Methods and Results} % Start the appendix part

\appendix

\setcounter{section}{0} \setcounter{table}{0} \setcounter{figure}{0}
\setcounter{equation}{0}
\setcounter{page}{1}

\renewcommand{\thesection}{S.\arabic{section}}
\renewcommand{\thetable}{S.\arabic{table}}
\renewcommand{\thefigure}{S.\arabic{figure}}
\renewcommand{\theequation}{S.\arabic{equation}}


Supplementary Material for `Accounting for the bin structure of data when
  fitting size spectra' by Andrew M.~Edwards, James P. W. Robinson, Julia
  L. Blanchard, Julia K. Baum and Michael J. Plank. Submitting to
  \emph{Marine Ecology Progress Series}.\\


% \addcontentsline{toc}{section}{Appendix} % Add the appendix text to the document TOC
\parttoc % Insert the appendix TOC

% \tableofcontents



\section{Extended methods and results}

\subsection{Outline}

In Section~\ref{sec:ibts} we give further details regarding the IBTS data.
In Section~\ref{sec:eight} we analyse the IBTS data by assuming the
`Length class' value to be the
exact length of the fish (as is often done, e.g.~\citealt{blan05, dgpr05}),
ignoring the fact that a length class of, say, 35~cm actually represents fish in
the range 35-36~cm. We explain how to extend each of the eight methods for
estimating the size-spectrum exponent to deal with the non-integer counts of the
numbers of fish. For the likelihood method we derive the likelihood function
for counts data (Section~\ref{sec:likeCount}) and show that it is valid for
non-integer counts (Section~\ref{sec:likeNonInt}).

We then derive the likelihood function for the MLEbins method, which properly accounts
for the species-specific body-mass bins (Section~\ref{sec:MLEbinslike}).
The other methods cannot be extended in a similar way.

In Section~\ref{sec:bodyMassBins} we show the body-mass bins for the IBTS data for
the species not shown in Figure~\ref{fig:nSeaFung/nSeaFungBins1}. In
Section~\ref{sec:ISD} we develop a method for plotting the data and resulting
ISD in a way that accounts for the bin structure and the uncertainty in the
estimate of the exponent $b$, and show figures for each year of the IBTS data.

Section~\ref{sec:furtherSim} gives summary statistics for the main simulation
results, plus sensitivity results from setting a larger minimum body size
and not sampling the smallest organisms. In Section~\ref{sec:hists} we
show histograms of a simulated data set from a bounded power-law distribution
to illustrate the highly skewed nature of the distribution.

We give details of how to access the freely available R code in
Section~\ref{sec:code}, in order to reproduce all results in this manuscript and
apply the methods to other data.

\subsection{Details of the IBTS data\label{sec:ibts}}

We obtained data collected by the North Sea IBTS from the ICES online Database
of Trawl Surveys ({\tt
  http://www.ices.dk/marine-data/data-portals/Pages/DATRAS.aspx}). Seven
countries survey the whole of the North Sea, bounded by the eastern English
Channel, northern continental margin above the Shetland archipelago, the east
coast of the UK and the Skagerrak and Kattegat regions \citep{ices15}. Surveys
take place in the first and third quarters of each year following standardized
gear and data collection protocols \citep{ices15}. We downloaded the number of
organisms caught per hour (surveyed catch per unit effort) for each species and
length class for all surveys in Areas~1 to~7 carried out in Quarter~1
from 1986-2015. Following
the protocol of \citet{ffrr12}, we restricted the time coverage of the surveys
to 1986 onwards because these are the years in which standardized fishing gear
(the GOV trawl -- chalut \`a grande ouverture verticale) was deployed on all
vessels. We downloaded data as ``cpue per length per haul'' in the formats
``Exchange Data'' and ``SMALK'', and applied the classifications published in
\citet{ffrr12} to remove all non-demersal fish.
During data collection on surveys, lengths were recorded as rounded
down to the nearest 1~cm [or 0.5~cm for Atlantic Herring (\emph{Clupea
  harengus}) and European Sprat (\emph{Sprattus sprattus})].
%Further
%processing was done to remove un-needed columns, remove counts of zero fish,
%convert lengths from mm to cm, and rearrange rows into a logical order.
We averaged counts of the same year/species/length-class combination
across the seven areas.
For simplicity this ignores potential differences in
sampling effort due to variations in haul duration, trawl speed and net area,
and fits a common size spectrum for the whole region.



\subsection{Applying eight methods to the IBTS data\label{sec:eight}}

We use eight methods that have been previously used to calculate size-spectra
slopes or exponents. See \citet{\MEE} for full details and example references for each method.
Briefly, the methods are:
\begin{itemize}
  \item Llin (Log-linear transform) -- plot linearly binned data on
log-linear axes then fit regression of log(count in bin) against midpoint of
bin;
%. & \citet{dgpr05}\\
  \item LT (log transform) -- plot linearly binned data on log-log axes then fit
regression of log(count in bin) against log(midpoint of bin);
%. & \citet{rg96}, \citet{bblhw12}\\
%
\item LTplus1 (log transform plus 1) -- plot linearly binned data on
$\logten$-$\logten$ axes then fit regression of $\logten (\mbox{count}+1)$
against $\logten (\mbox{midpoint of bin})$;
%. & \citet{dpmg04}, \citet{gdjp05}\\
%
\item LBmiz (logarithmic binning as done in the R package {\tt mizer} by
\citealt{sba14}) -- bin data using $\logten$ bins, but with largest bin the same
arithmetic size as the penultimate bin, then fit regression of $\log$(count in
bin) against $\log$(lower bound of bin);
%. & Scott et al.'s (2014) {\tt mizer} R package\\
%
\item LBbiom (logarithmic binning) and then fit biomass size spectrum -- bin body
masses using $\logtwo$ bins then fit regression of
$\logten (\mbox{biomass in bin})$ against $\logten (\mbox{midpoint of bin})$;
% . & \citet{mj06}, \citet{jdw07}, \citet{tdss15}\\
%
\item LBNbiom (logarithmic binning with normalisation and then fit biomass size
spectrum) -- bin body masses using $\logtwo$ bins, then fit regression of
$\logten$~(biomass in bin divided by bin width) against $\logten
(\mbox{midpoint of bin})$;
%. & \citet{blan05}, \citet{rps11}\\
%
\item LCD (logarithmic plotting of $1 - F(x)$; i.e.~one minus the cumulative
distribution) -- rank data from largest (rank~1) to smallest (rank~$n$) and fit
regression of $\log(\mbox{rank}(x) / n)$ against $\log~x$;
%. & \citet{vpbd97}, \citet{rbm14}\\
%
\item MLE (maximum likelihood estimate) -- calculate maximum likelihood estimate of
$b$.
% Data and fitted distribution can be plotted on a rank/frequency plot. & \citet{abbzza11}, \citet{rb16} \\
\end{itemize}


Using each method we calculate the size-spectra exponent $b$ separately for
each year
of data.
As in Table~\ref{tab:egData} (following \citealt{ffrr12}) we use the
`Length class' values (the minimum of each length bin) to calculate body masses
associated with the numbers of individuals.
Although some methods use counts of individuals and some use biomass,
they are related through (\ref{PLB}) and (\ref{biomass}) and hence can estimate
$b$ \citep{\MEE}. This gives us a time series of estimated annual values of $b$
with confidence intervals, one time series for each method. We then fit a
weighted linear regression (that uses the confidence intervals -- see below)
to each time series to see whether there is a
significant change ($p < 0.05$) in $b$ through time.

However, the eight methods first have to be extended to deal with
the non-integer counts of numbers of fish of each body mass. For the
binning-based methods (Llin, LT, LTplus1, LBmiz, LBbiom and LBNbiom) the
extension is fairly obvious. The bins used to fit the size spectra are defined
by the method.
% , and the total number of individuals in a bin is a simple sum of the numbers of individuals that have a body mass that lies within that bin (which can be done computationally as a histogram calculation in \texttt{R}).
Each count of a particular body mass (each row in Table~\ref{tab:egData}) is
assigned to a bin that is defined by the method being used, and (for each year)
the total count in each body-mass bin is calculated by summing the `Number'
values for that bin. The resulting total count in each body-mass bin no longer
has to be an integer (because the `Number' values are not integers), but the
resulting regression fits can be calculated as described by \citet{\MEE}.

The LCD method requires ranking the body mass, $x$, of each individual fish in
descending order from 1 (largest) to $n$ (smallest) and then fitting a
regression of log(fraction of values $\geq x$) against log~$x$. However, for the
IBTS data the body masses do not correspond to individual fish, rather we have
the number of fish of a certain body mass -- but that number can be
non-integer. We use the simplest approach to adapt the LCD method. First, for a
particular year arrange the data (`Number' for each body mass from
Table~\ref{tab:egData}) in descending order of body mass, and for each row
calculate the cumulative sum of the `Number' -- this is akin to assigning ranks
for individual body masses. Dividing the cumulative sum by the total `Number',
gives the fraction of individuals with a body mass $\geq x$ for each given body
mass, $x$. Then plot log(proportion of values $\geq x$) against log~$x$ and fit
a regression, as per Figure~2(g) of \citet{\MEE}. Each plotted point no longer
corresponds to an individual fish, but relates to a row in
Table~\ref{tab:egData}.

The original MLE method of \citet{\MEE} requires the
body mass of individual fish, and so in Sections~\ref{sec:likeCount}
and~\ref{sec:likeNonInt}
we formally extend it to deal with count data.

So we calculate the estimated $b$ for each of the 30~years of data, using each
of the eight methods. For each method, we plot the resulting time series of
estimates of $b$ and analyse any trend by fitting a weighted linear regression
to the estimates. A weighted regression is used because we can estimate the
variances as the square of the standard errors from regression outputs, and the
variances appear to vary over time. For the MLE method, having already
calculated the 95\% confidence intervals $(b_{\text{low}}, b_{\text{high}})$
using the profile likelihood method
\cite{hm97}, the standard error, $\gamma$, can be calculated from the
confidence interval being approximately $b_{\text{MLE}} \pm 1.96 \gamma$ where
$b_{\text{MLE}}$ is the MLE for $b$ (Burnham and Anderson, 2002).
% Manually since not in main text
Explicitly, we used the mean
\eb
\gamma = \frac{ |b_{\text{low}} - b_{\text{MLE}} | +
  |b_{\text{high}} - b_{\text{MLE}} | }{2 \times 1.96}
  = \frac{ b_{\text{high}} - b_{\text{low}} }{2 \times 1.96}.
\ee
% p174 ba02.

Figure~\ref{fig:nSeaFung/nSeaFungTrends} and Table~\ref{tab:trendRes}
% From fitting-size-spectra/ nSeaFungAnalysis.Snw .tex results and from
%  nSeaMLEbins.Snw .tex for MLEbins
% latex table generated in R 3.4.3 by xtable 1.8-2 package
% Tue May 15 15:12:04 2018
\begin{table}[tp]
\centering
\caption{Results of weighted regression analyses of trend through time of the
  estimated exponent $b$ for the IBTS data,
  as estimated using each of the eight methods shown
  in Figure~\ref{fig:nSeaFung/nSeaFungTrends}, plus the MLEbins method
  developed here.
  `Trend' is the estimated annual
  trend, with 95\% confidence intervals given by `Low' and `High', $p$ is the
  $p$-value for the probability that the trend is significantly different to 0,
  and $R^2$ is the coefficient of determination. If $p \geq 0.05$ then the
  trend can be considered not significantly different to 0. If $p<0.05$ then a
  negative trend indicates a statistically significant decline in the exponent
  over time, and a positive trend indicates a statistically significant
  increase.}
\label{tab:trendRes}
\begin{tabular}{lrrrrr}
\hline Method & Low & Trend & High & $p$ & $R^2$ \\    % CI width
\hline
Llin & -0.0000 & -0.0000 & -0.0000 & 0.01 & 0.19 \\    %<0.0001
LT & -0.0313 & 0.0052 & 0.0417 & 0.77 & 0.00 \\        % 0.0104
LTplus1 & -0.0206 & 0.0058 & 0.0321 & 0.66 & 0.01 \\   % 0.0115
LBmiz & -0.0060 & -0.0034 & -0.0009 & 0.01 & 0.21 \\   % 0.0051
LBbiom & -0.0076 & -0.0042 & -0.0008 & 0.02 & 0.19 \\  % 0.0068
LBNbiom & -0.0076 & -0.0042 & -0.0008 & 0.02 & 0.19 \\ % 0.0068
LCD & -0.0057 & -0.0030 & -0.0002 & 0.04 & 0.15 \\     % 0.0055
MLE & -0.0047 & -0.0010 & 0.0027 & 0.60 & 0.01 \\      % 0.0020
\hline
MLEbins & -0.0043 & -0.0010 & 0.0024 & 0.56 & 0.01 \\  % 0.0019
\hline
\end{tabular}
\end{table}
shows that the absolute estimates of
$b$ are highly dependent upon the method used. The confidence intervals are
somewhat narrower ($\leq$0.002) in absolute size for the Llin and MLE methods than those for
the other methods ($\geq$0.005),
and vary in size through time. There is no significant change in $b$ when
using three of the estimation methods (namely LT, LTplus1 and MLE), yet a
significant negative decline when using the remaining five methods. Thus, five
methods imply a steepening of the size spectrum over time, whereas three methods
imply no change. This demonstrates how methodological differences can lead to
differing ecological conclusions.  % check we've mentioned steepening already

\subsection{Likelihood function for count data\label{sec:likeCount}}

Here we derive the likelihood function for count data, that is then shown in
Section~\ref{sec:likeNonInt} to apply to non-integer counts. The latter was used for
the MLE method (that ignores the bin structure) for the IBTS data.

From \citet{\MEE}, the log-likelihood function for the PLB model for $b \neq -1$
is
\eb
\log [ \mbox{L} (b | \mbox{data}~{\bf x})] & = & \sum_{j=1}^n \log
f(x_j)\\ & = & n \log\left( \dfrac{ b + 1}{\xmax^{~b + 1} - \xmin^{~b+1}}
\right) + b \sum_{j=1}^n \log x_j,
\label{llPLB}
\ee
where L$(b | \mbox{data}~{\bf x})$ is the likelihood of a particular value
of the unknown parameter $b$ given the known data ${\bf x} = \{x_1, x_2, x_3,
..., x_n\}$, and $f(\cdot)$ is the probability density function (\ref{PLB}). The
maximum likelihood estimates for $\xmin$ and $\xmax$ are simply the minimum and
maximum values of the data.

Now, if some of the $x_j$ are repeated values, we can represent the data in a
more concise form. For example, the data set ${\bf x} = \{4, 6, 6, 6, 9, 10,
10\}$ can be represented as counts $\{1, 3, 1, 2\}$ for each of the unique $x$
values $\{4, 6, 9, 10\}$. More generally, let $c_k$ be the count (number of
repetitions) of $x'_k$, where $k=1, 2, 3, ..., K \leq n$, and $\sum_{k=1}^K c_k
= n$. The ${\bf x'} = \{x'_k\}$ represent the unique values of the original
$x_j$, and $K=n$ only when all $c_k = 1$ (i.e.~the original $x_j$ values are all
unique).

To obtain the log-likelihood function, note that \eb \sum_{j=1}^n \log f(x_j) =
\sum_{k=1}^K c_k \log f(x'_k),
\label{countsEquiv}
\ee because if $x_j$ is repeated $c_k$ times then the corresponding contribution
of the $x_j$ values to the overall $\sum \log f(x_j)$ is just $c_k \log f(x'_k)$
with $x'_k = x_j$. Hence, the equivalent log-likelihood function to
(\ref{llPLB}) is \eb \log [ \mbox{L} (b | \mbox{data}~{\bf x'})] & = &
\sum_{k=1}^K c_k \log f(x'_k)\\ & = & n \log\left( \dfrac{ b + 1}{\xmax^{~b + 1}
  - \xmin^{~b+1}} \right) + b \sum_{k=1}^K c_k \log x'_k.
\label{llPLBcount}
\ee

Similarly, for $b = -1$ the log-likelihood function is \eb \log [ \mbox{L} (b =
  -1 | \mbox{data}~{\bf x'})] = - n \log \left( \log \xmax - \log \xmin \right)
- \sum_{k=1}^K c_k \log x'_k.
\label{llPLB2count}
\ee

Note that if a value $x_j$ is repeated because it represents a value that is
rounded due to the resolution of the measurement process, (e.g.~$x_j = 10$~g
really represents a value in the range $9.5-10.5$~g and then occurs multiple
times in the data set) then the MLEbin method should be used \citep{\MEE}. If
the values $\{ x_j \}$ each represent true discrete measurements, then the
\emph{continuous} power-law distribution (\ref{PLB}) is not appropriate and a
\emph{discrete} distribution should be used. But for
length and weight measurements used for size spectra, the underlying variables
are indeed continuous, and any discrete values are due to measurement
resolution.

\subsection{Likelihood function for data with non-integer counts\label{sec:likeNonInt}}

% **The reason for suggesting this is so the likelihood approach can be applied to situations where raw count data get scaled to give non-integer values. \marginpar{Okay Julia Bl.?} For example, \citet{blan05} analysed data from annual groundfish surveys in the Celtic Sea. Since the duration of individual trawls was not constant, the integer counts of the number of fish within each size class in a particular trawl needed to be divided by that trawl's duration, to give the number of fish per hour within each size class. This standardised the unit of measurement across trawls of different duration. The resulting counts per hour, $d_j$, are therefore not restricted to being integers. Similarly, \citet{dgpr05} analysed bottom-trawl survey data that were `catch in number by size or Lmax [maximum length] class per hour fishing'; these can also result in non-integer counts. We ascertain that the above likelihood functions can still be applied. This provides a formal statistical method for dealing with such data.

For the IBTS data, the counts $c_k$ corresponding to each body mass $x'_k$ are
not all integers -- they can take non-integer values. The duration of individual
research trawls was not constant. So the integer counts of the number of a
particular species of fish within a particular length bin from a particular
trawl needed to be divided by that trawl's duration, to give the number of fish
per hour within each length bin. This standardised the unit of measurement
across trawls of different duration. The resulting counts per hour, $c_k$, are
therefore not restricted to being integers. Similarly, \citet{dgpr05} analysed
bottom-trawl survey data that were `catch in number by size or Lmax [maximum
  length] class per hour fishing' that resulted in non-integer counts, as did
\citet{blan05} in their analysis of groundfish surveys from the Celtic Sea.

We ascertain that the above likelihood functions (\ref{llPLBcount}) and
(\ref{llPLB2count}) can still be applied to data containing non-integer
counts. In the above derivations, each count $c_k$ of the number of fish of size
$x'_k$ was assumed to be integer-valued. But we now show that this condition can
be relaxed, and therefore allow each $c_k$ to be non-integer.
If $c_k$ are non-integer then they can be scaled by a constant factor $z$ to
give counts $c'_k=z c_k$ such that all the $c'_k$ are integers. Using $c'_k$ in
place of $c_k$ and $n'=z n$ in place of $n$ simply multiplies the log-likelihood
function in (\ref{llPLBcount}) by a constant $z$, which does not change the
location of
its maximum. Hence, using non-integer densities $c_k$ in (\ref{llPLBcount})
gives the same MLE for $b$ as scaling the densities up to integer values.

However, any such rescaling, by multiplying (\ref{llPLBcount}) by $z$,
changes the curvature of the log-likelihood function, and so changes the
width of the resulting confidence interval. This is because the interval
is calculated using the profile likelihood-ratio test as the
range of parameters for which the log-likelihood is within 1.92 of the maximum
value of the log-likelihood (page~163 of \citealt{hm97}). The 1.92 is based on
a chi-squared distribution with one degree of freedom (and is more precisely
calculated as {\tt qchisq(0.95,1)/2} in our R code), and is an absolute value
that does not change when the log-likelihood function is scaled by $z$.

Thus, calculated confidence intervals will depend on the units of the counts --
changing the counts of fish in the IBTS data from numbers per hour of trawling
to numbers per day of trawling will change the confidence intervals of $b$ (but
not the MLEs). An analogy is tossing 10 identical (but unfair)
coins and getting 7 heads and 3 tails
(the MLE for the probability of getting a head is 0.7, but with a wide
confidence interval), compared to tossing 1,000 coins and getting 700 heads and
300 tails (the MLE is still 0.7, but the uncertainty will be much narrower given
the larger sample size).
%[**much less likely to include 0.5 (fair coin) - could
%calculate some numbers].

Note that the units of measurement (g vs.~kg) will not affect
the MLEs or the confidence intervals as this change only adds a constant value
to the log-likelihood function.

\subsection{Likelihood function for the MLEbins method\label{sec:MLEbinslike}}

Here we derive the likelihood function for the MLEbins method. This extends the
MLEbin method for when the body-mass bins are not defined the same for all species. This
occurs for the IBTS data because the counts in the
1-cm (or 0.5-cm) length bins yield counts in body-mass bins via length-weight
relationships. Since the length-weight relationships are species-specific,
the resulting body-mass bins are species-specific, as shown in
Figures~\ref{fig:nSeaFung/nSeaFungBins1} and
\ref{fig:nSeaFung/nSeaFungBins2}-\ref{fig:nSeaFung/nSeaFungBins4}.

We extend and generalise the MLEbin method derived in \citet{\MEE}, which in
turn extended the methods from Edwards et~al.~(2007)
% manually since not in main text
and \citet{edwa11} for binned
movement data. The aim is to obtain the likelihood functions to calculate the
maximum likelihood estimate for the exponent $b$ in (\ref{PLB}). The MLEbin
method in \citet{\MEE} assumed that all data were binned using the same binning
protocol, such that there is just one set of bin breaks that gives the
breakpoints between the bins. Here we relax that assumption to allow for
species-specific bin breaks, as occurs when converting binned length data into
binned body-mass data using species-specific length-weight relationships.

Consider the data to consist of counts, $d_{sj}$, of the number of individuals
of species $s$ that are in each size bin ${j = 1,2,3,...,J_s}$, where $J_s$ is
the index of the final bin for species $s$, and there are $S$ species labelled
$s=1, 2, 3, ..., S$.

Let bin $j$ cover the values of $x$ (which could be weight, or length for
fitting a length size spectrum) in the
interval $[w_{sj}, w_{s, j+1})$, such that $w_{s1}, w_{s2}, ..., w_{s, J_s +1}$
  define the bin breaks. For example, for species $s=3$, bin $j=5$ goes from
  $w_{3,5}$ to $w_{3,6}$. [In terms such as $w_{sj}$ the $s$ and $j$ are
    indices; we include commas where necessary to avoid ambiguity, such as
    $w_{s, j+1}$ and $w_{3,6}$]. For bin $j=J_s$ the interval is $[w_{s J_s},
    w_{s, J_s +1}]$, which includes the upper bound. The sample size (total
  number of counts) for species $s$ is $n_s = \sum_{j=1}^{J_s} d_{sj}$, and the
  overall total number of counts is $n = \sum_{s=1}^{S} n_s$. We assume that,
  for each species, the first and last bins each have at least one individual in
  them (i.e. $d_{s,1}, d_{s, J_s} > 0$).

Similar calculations to those by Edwards et~al.~(2012)
% manually since not in main text
for unbinned data show that the
known $\min \{w_{s1}\}_{s=1}^S$ and $\max \{w_{s, J_s +1}\}_{s=1}^S$ are the
maximum likelihood estimates for $\xmin$ and $\xmax$, respectively. So $\xmin$
and $\xmax$ can be set to their respective maximum likelihood estimates, and we
now only need to calculate the maximum likelihood estimate of $b$ to fully
specify the ISD (\ref{PLB}).

The probability that the body size of an individual lies within the range
$[w_{sj}, w_{s, j+1})$ is simply (assume for now that $b \neq -1)$ \eb
  \mbox{P}\left(\mbox{individual has body size in }[w_{sj}, w_{s, j+1}) |
    b\right) & = & \int_{w_{sj}}^{w_{s, j+1}} C x^{b} \mbox{d}x\\ & = &
    \frac{C}{b + 1} \left[ x^{b + 1}\right]_{w_{sj}}^{w_{s, j+1}}\\ & = &
    \frac{C}{b + 1} \left( w_{s, j+1}^{b + 1} - w_{sj}^{b + 1} \right),~~~~\\ &
    = & \frac{\cancel{(b + 1)}\left( w_{s, j+1}^{b + 1} - w_{sj}^{b + 1}
      \right)}{\cancel{(b + 1)}(\xmax^{b+1}-\xmin^{b+1})},\\ & = & \frac{w_{s,
        j+1}^{b + 1} - w_{sj}^{b + 1}}{\xmax^{b+1}-\xmin^{b+1}}.
\label{Pind}
\ee The individual can be from any species (not necessarily species $s$).

Recall that the $w_{sj}$ are all known, since they are the species-specific bin
breaks. Thus the only unknown in (\ref{Pind}) is $b$. The remaining known
quantities are the counts, $d_{sj}$, of the number of individuals of species $s$
that are in each bin $[w_{sj}, w_{s, j+1})$.

Given these counts, we develop a multinomial log-likelihood function
\citep{lawl03} as follows. The log-likelihood function for the parameter $b$
(the only unknown), given the counts $\{ d_{sj} \}_{j=1}^{J_s}$ corresponding to
bin breaks $\{ w_{sj} \}_{j=1}^{J_s +1}$ for each species $s=1, 2, 3, ..., S$,
is
\eb
l(b | \{d_{sj} \}, \{w_{sj} \}) & = & \log \left[ \prod_{s=1}^{S}
  \prod_{j=1}^{J_s} \left( \mbox{P}(\mbox{individual has body size in
  }[w_{sj}, w_{s, j+1}) | b) {\vrule height 12pt width 0pt} \right)^{d_{sj}}
    \right]
~~~~~~~~~~~\label{ll1}\\ % \vrule makes () larger
& = & \sum_{s=1}^{S} \sum_{j=1}^{J_s} d_{sj} \log \left(
  \mbox{P}(\mbox{individual has body size in }[w_{sj}, w_{s, j+1}) | b) {\vrule
      height 12pt width 0pt} \right)
~~~~~~~~~~~\label{ll2}\\
& = & \sum_{s=1}^{S} \sum_{j=1}^{J_s} d_{sj} \log \left( \frac{w_{s, j+1}^{b +
        1} - w_{sj}^{b + 1}}{\xmax^{b + 1} - \xmin^{b + 1}}
    \right)
\label{abs1}\\
& = & \sum_{s=1}^{S} \sum_{j=1}^{J_s} d_{sj} \left(
    \log \left| w_{s, j+1}^{b + 1} - w_{sj}^{b + 1}\right| - \log \left|
    \xmax^{b + 1} - \xmin^{b + 1} \right| {\vrule height 14pt width 0pt}
    \right)
\label{abs2}\\
%
%
%& = & \sum_{s=1}^{S} \left[ \sum_{j=1}^{J_s} d_{sj} \log \left| w_{s, j+1}^{b + 1} - w_{sj}^{b + 1}\right| - \sum_{j=1}^{J_s} d_{sj} \log \left| w_{s, J_s +1}^{b + 1} - w_{s1}^{b + 1} \right| \right]\\
%& = & \sum_{s=1}^{S} \left[ \sum_{j=1}^{J_s} d_{sj} \log \left| w_{s, j+1}^{b + 1} - w_{sj}^{b + 1}\right| - \log \left| w_{s, J_s +1}^{b + 1} - w_{s1}^{b + 1} \right| \sum_{j=1}^{J_s} d_{sj} \right]\\
%& = & \sum_{s=1}^{S} \left[ \sum_{j=1}^{J_s} d_{sj} \log \left| w_{s, j+1}^{b + 1} - w_{sj}^{b + 1}\right| - n_s \log \left| w_{s, J_s +1}^{b + 1} - w_{s1}^{b + 1} \right| \right]\\
& = & - \log \left| \xmax^{b + 1} - \xmin^{b + 1} \right| \sum_{s=1}^{S}
    \sum_{j=1}^{J_s} d_{sj} + \sum_{s=1}^{S} \sum_{j=1}^{J_s} d_{sj} \log
    \left| w_{s, j+1}^{b + 1} - w_{sj}^{b + 1}\right|\\
& = & - n \log \left|
    \xmax^{b + 1} - \xmin^{b + 1} \right| + \sum_{s=1}^{S} \sum_{j=1}^{J_s}
    d_{sj} \log \left| w_{s, j+1}^{b + 1} - w_{sj}^{b +
      1}\right|.
\label{loglikbin}
\ee
%
The two terms inside the absolute symbols $| \cdot |$, i.e. $w_{s, j+1}^{b + 1}
- w_{sj}^{b + 1}$ and $\xmax^{b + 1} - \xmin^{b + 1}$, are both positive for $b
< -1$ and both negative for $b>-1$ (because $w_{s, j+1} > w_{sj}$ and $\xmax >
\xmin$ by definition), such that taking their absolute values ensures that
(\ref{abs1}) and (\ref{abs2}) are equivalent. Equation (\ref{loglikbin}) is
analogous to (A.70) in the Appendix of \citet{\MEE} for which the data have a
single non-species-specific set of bin breaks, but extended here to account for
the species-specific bin breaks. It cannot be analytically solved to give the
maximum likelihood estimate of $b$ (by differentiating with respect to $b$ and
setting to 0), and so numerical methods are required. It is formulated in the
function {\tt negLL.PLB.bins.species()} in our R code.

Note that (\ref{ll1}) necessarily assumes that individual body sizes are
independently distributed, which may be an assumption that could be investigated
in later work. However, we have not had to make any assumptions about the
distribution of body sizes within each species, only about the community-level
ISD.

For the case where $b = -1$, we have $C = 1 / (\log \xmax - \log \xmin)$ from
(\ref{PLB.C}) and, analogous to (\ref{Pind}) we have
\eb
\mbox{P}\left(\mbox{individual has body size in }[w_{sj}, w_{s, j+1}) | b =
  -1\right) & = & \int_{w_{sj}}^{w_{s, j+1}} C x^{-1} \mbox{d}x\\ & = & C \left[
  \log x {\vrule height 12pt width 0pt} \right]_{w_{sj}}^{w_{s, j+1}}\\
  & = & \frac{\log w_{s, j+1} - \log w_{sj}}{\log \xmax - \log \xmin}.
\label{Pindbmin1}
\ee

Analogous to (\ref{ll2}), the log-likelihood function is then just
\eb
\nonumber l(b =-1| \{d_{sj} \}, \{w_{sj} \}) & = & \sum_{s=1}^{S} \sum_{j=1}^{J_s}
d_{sj} \log \left( \mbox{P}\left(\mbox{individual has body size in }[w_{sj},
    w_{s, j+1}) | b = -1\right) {\vrule height 12pt width 0pt} \right)\\
& & \\
& = & \sum_{s=1}^{S} \sum_{j=1}^{J_s} d_{sj} \log \left( \frac{\log w_{s, j+1}
    - \log w_{sj}}{\log \xmax - \log \xmin}\right)\\
\nonumber & = &
  \sum_{s=1}^{S} \left[ \sum_{j=1}^{J_s} d_{sj} \left( \log \left( \log w_{s,
      j+1} - \log w_{sj}\right) - \log \left( \log \xmax - \log \xmin \right)
  {\vrule height 14pt width 0pt} \right) \right]\\
& & \\
\nonumber & = & -
  \log \left( \log \xmax - \log \xmin \right) \sum_{s=1}^{S} \sum_{j=1}^{J_s}
  d_{sj} + \sum_{s=1}^{S} \sum_{j=1}^{J_s} d_{sj} \log \left( \log w_{s, j+1}
    - \log w_{sj}\right)\\
  & & \\
\nonumber & = & - n \log \left( \log \xmax -
  \log \xmin \right) + \sum_{s=1}^{S} \sum_{j=1}^{J_s} d_{sj} \log \left( \log
  w_{s, j+1} - \log w_{sj}\right).
\label{llbinbmin1}\\
\ee

Note that the equivalent single-species equation (A.75) in \citet{\MEE} was
incorrect~-- it contained the $\log$ terms that results from integrating
$x^{-1}$, but not the $\log$ term resulting from taking the log-likelihood. The
correct equation is \eb l(b = -1 | \mbox{data}) & = & - n \log (\log w_{J +1} -
\log w_{1}) + \sum_{j=1}^{J} d_{j} \log ( \log w_{j+1} - \log
w_{j}).\label{llMEEcorrected} \ee This is equivalent to (\ref{llbinbmin1}) with
$S=1$ and dropping the $s$ subscript (and replacing $\xmin$ and $\xmax$ by their
maximum likelihood estimates of $w_{1}$ and $w_{J+1}$, respectively). This error
has now been corrected in the {\tt negLL.PLB.binned()} function in {\tt
  PLBfunctions.r} in the GitHub code repository associated with \citet{\MEE},
and is documented there as Issue~7 at {\protect{\tt
  https://github.com/andrew-edwards/fitting-size-spectra/issues/7}} --
thank you to Philip Wallhead (Norwegian Institute for Water Research,
pers.~comm.) who
independently noticed this error and that we had not originally fully corrected it.
In practice
this error is very unlikely to matter, as it will only occur when the
log-likelihood function {\tt negLL.PLB.binned()} is called for a value of
precisely $b=-1$, which is unlikely to occur (due to machine precision) when,
for example, minimising the negative log-likelihood function. This omission did
not occur for the unbinned data (equation (A.11) in \citealt{\MEE}), as can be
seen by the {\tt log( log(xmax) - log(xmin) )} term in the function {\tt
  negLL.PLB()} in {\tt PLBfunctions.r}.

\subsubsection{Non-integer counts}

We have assumed in (\ref{ll1}) that the counts $\{d_{sj}\}$ are integer
valued. However, using the same argument developed for (\ref{countsEquiv}) and
(\ref{llPLBcount}) we contend that this assumption can be relaxed, i.e.~that the
$\{d_{sj}\}$ can take non-integer values. This is what we have for the IBTS
data, since the counts are counts per hour of trawling (so that the data are
standardised). For example, if three individual Atlantic Cod whose size falls in
bin $j$ are caught during a trawl of 17 minutes, then the standardisation yields
a count $d_{sj} = 3 \times 17/60 = 0.85$ which is non-integer. Averaging counts
per hour across areas also generates non-integer counts.

Our approach (and R code) is applicable to any data
set that has species-specific bin breaks. It can be used directly for length
data or for body-mass data. For the IBTS data, calculations of the individual
size distribution based on lengths (the length size spectrum), rather than body
masses, would still require this approach because two of the species (herring and
sprat) used 0.5-cm bins, compared to 1-cm bins for all the others. In practice,
to simplify the numerical computations, the analysis could just consider two
pseudo-species, $s=1$ representing all species for which 1-cm bins were used
(and the counts within each length bin just summed across all such species), and
$s=2$ for the 0.5-cm species (herring and sprat).

\subsection{Resulting body-mass bins for IBTS data\label{sec:bodyMassBins}}

Figures~\ref{fig:nSeaFung/nSeaFungBins2}-\ref{fig:nSeaFung/nSeaFungBins3} show
the equivalent of Figure~\ref{fig:nSeaFung/nSeaFungBins1} for the remaining 90
species of the IBTS data (and Figure~\ref{fig:nSeaFung/nSeaFungBins4} shows an
enlargement of part of Figure~\ref{fig:nSeaFung/nSeaFungBins3}).

The body-mass bin with the biggest ratio of its width to its lower bound
occurs for the Blackbelly Rosefish (\emph{Helicolenus dactylopterus}) which is
species code 127251, indicated by $\times$ in
Figure~\ref{fig:nSeaFung/nSeaFungBins4}. The bin goes from 10.29~g to 20.31~g,
such that the ratio of bin-width to the lower bound is 0.97 (though this is hard
to see in Figure~\ref{fig:nSeaFung/nSeaFungBins4}). Consequently, in our
earlier eight-methods
analysis that did not account for the bin structure, all individuals of this
species with a body-length of 4-5~cm would get
assigned a body length of 4~cm, and consequently a body mass of
10.29~g. However, the true possible range of body masses is 10.29-20.31~g. Thus,
individuals at the upper end of this range would be assigned body masses of
about half of their actual body mass.

A similar effect occurs for \emph{all} body masses in the data set. This
systematic rounding down of body masses impacts all the methods except for the
MLEbins method that explicitly accounts for the bin structure.

The widest resulting body-mass bin occurs for Atlantic Cod (\emph{Gadus
  morhua}), which is species code 126436, the rightmost species in
Figure~\ref{fig:nSeaFung/nSeaFungBins3}. The widest bin goes from 35.630~kg to
36.462~kg, with a width of 832~g.

\onefig{nSeaFung/nSeaFungBins2}{As for Figure~\ref{fig:nSeaFung/nSeaFungBins1}
  but for the next 45 species -- note the different vertical axis scale from
  Figure~\ref{fig:nSeaFung/nSeaFungBins1}.}
% Shorthorn Sculpin is 127203, but maybe not worth saying anything.

\onefig{nSeaFung/nSeaFungBins3}{As for Figure~\ref{fig:nSeaFung/nSeaFungBins1}
  but for the final 45 species -- note the different vertical axis scale from
  Figure~\ref{fig:nSeaFung/nSeaFungBins1}. If the figure is rotated
  90$^\circ$ clockwise it is somewhat characteristic of a bounded
  power-law distribution, with a very few species growing much larger than the
  remaining species.}

\onefig{nSeaFung/nSeaFungBins4}{As for Figure~\ref{fig:nSeaFung/nSeaFungBins3}
  but just for species with maximum body mass up to 10~kg to more clearly show
  their body-mass bins. The $\times$ for species code 127251 indicates
  Blackbelly Rosefish (\emph{Helicolenus dactylopterus}) which has the
  body-mass bin with the largest ratio of width to lower bound (see text).}

\clearpage

% From nSeaFungMLEbinsRecommend-ISD.Snw
\subsection{Plotting the IBTS data and the resulting PLB fit for
  each year\label{sec:ISD}}

Having accounted for the species-specific body-mass bins when fitting
the ISD to the IBTS data, we now
describe how to plot the data (and resulting fit) in a way that also accounts
for the binning.
In Figure~6 of \citet{\MEE} we recommended plotting both the biomass size
spectrum (in binned form) and the abundance size spectrum.
However, since the biomass size spectrum requires further binning
(which is problematic, as shown in Figure~\ref{fig:binFigs/binLWlog})
here we focus on the abundance size spectrum.

For unbinned data, in \citet{\MEE} we suggested plotting the abundance size
spectrum by ranking the body mass, $x$, of each individual fish in
descending order from 1 (largest) to $n$ (smallest), and plotting
log(number of values $\geq x$) against log~$x$ for each individual body mass $x$.
The fitted ISD is then plotted as $(1 - F(x))n$, where $F(x)$ is the cumulative
distribution function
\eb
F(x) = \left\{
          \begin{array}{ll}
            \dfrac{x^{b + 1} - \xmin^{~b + 1}}{\xmax^{~b + 1} - \xmin^{~b + 1}},
            \vspace{2mm}
            & b \neq -1\\
            \dfrac{\log (x / \xmin)}{\log (\xmax / \xmin)},
            & b = -1,
          \end{array}
      \right.
\label{Fx}
\ee
and $n$ is the total sample size.

For the binned IBTS data, we first try defining the body sizes using the minimum
body mass for each of the bins. For bin $j$ for species $s$ the minimum body mass
is, by definition, $w_{sj}$.
We use $D_{sj}$ to represent the equivalent to the unbinned
`number of values $\geq x$' ,
and define it as the total number of counts in all the bins that have minimum
body mass $\geq w_{sj}$:
\eb
D_{sj} = \sum_{s'j' {\mbox{~with~}} w_{s'j'} \geq w_{sj}} d_{s'j'}.
\label{Dsj}
\ee
The summation is over combinations of $s'$ and $j'$ that satisfy
$w_{s'j'} \geq w_{sj}$.
% Computationally, it helps to rank the bins.
We can then account for the
range of body masses in each bin by plotting a horizontal line from $w_{sj}$ to
$w_{s, j+1}$, i.e. showing the range of each bin, on the x-axis, and
$D_{sj}$ on the y-axis; these are shown as the green horizontal bars in
Figure~\ref{fig:ISD-1999.main} for the 1999 data.

However, we have overlapping bins (because of the species-specific
length-weight relationships), but have calculated $D_{sj}$ based only on the
ranking of the minima of each bin. Only using the minima implies, for example, that
a bin with range 20-22~g would be designated as containing individuals
larger than a bin with range 19-25~g, even though the latter bin can contain
individuals larger (e.g.~of mass 24~g) than the former bin. And bins may overlap without one
fully encompassing the other, such that using the minima means that a bin with range 30-35~g is considered
to always contain individuals larger than a bin with range 29-34~g, even though
the latter bin can again contain individuals larger than some in the former bin.

We account for such possibilities by calculating, for bin $sj$, the range
$[E_{sj1}, E_{sj2}]$, where
\eb
E_{sj1} & = & \sum_{s'j' {\mbox{~with~}} w_{s'j'} \geq w_{s,j+1}} d_{s'j'},
\label{Esj1}\\
\nonumber~ \\
E_{sj2} & = & \sum_{s'j' {\mbox{~with~}} w_{s',j'+1} > w_{sj}} d_{s'j'}.
\label{Esj2}
\ee
The low end of the range, $E_{sj1}$, accounts for bins whose minimum value
equals or exceeds
the maximum value of bin $sj$, i.e.~bins with $s'$ and $j'$ such that
$w_{s'j'} \geq w_{s,j+1}$, because individuals in such bins
must be larger than (or equal to)
those in bin $sj$ as the bins do not overlap
(see Figure~\ref{fig:binCountRangeSchematic}, further explained below).
%
\begin{figure}[tp]
\begin{center}
% \epsfxsize=7in
\epsfbox{../../fitting-size-spectra/codeBinStructure/nSeaFung/binCountRangeSchematic.eps}
\end{center}
\caption{Schematic diagram to explain how we calculate
  the range of counts of individuals that are larger than those in a given bin.
  Red and pink body-mass bins are those for Snakeblenny and Moustache Sculpin
  (here labelled species 1 and 2, respectively) from
  Figure~\ref{fig:nSeaFung/nSeaFungBins1}.
  Bin breaks are denoted by $w_{s,j}$
  and the number inside each bin is the number observed per hour of trawling.
  For illustration data are combined across all years and only bins with
  minima $>15$~g are shown. The target bin has $s=2$ and $j=7$ and therefore
  has bin breaks $w_{2,7}$ and $w_{2,8}$ and is indicated by the vertical grey
  lines.
  The first letter in each pair (`NN', `NY', or `YY') denotes whether or not
  each bin is included in the low count $E_{2,7,1}$, i.e. its lower bound is
  $\geq$ the upper bound of the target bin. Similarly, the second letter denotes
  whether or not each bin is included in the high count $E_{2,7,2}$, i.e.
  its upper bound is $>$ the lower bound of the target bin. Summing the
  respective counts as per (\ref{Esj1}) and (\ref{Esj2}) gives
  $E_{2,7,1} = 0.41$ and $E_{2,7,2} = 2.71$.
  }
  \label{fig:binCountRangeSchematic}
\end{figure}
%
The high end of the range, $E_{sj2}$, accounts for bins whose maximum value
exceeds the minimum value of bin $sj$, i.e.~bins with $s'$ and $j'$ such that
$w_{s',j'+1} > w_{sj}$, because individuals in such bins may be larger
than those in bin $sj$ as the bins overlap.
For $E_{sj2}$ we need to use the criteria of $>$, rather than $\geq$,
to avoid erroneously counting bins whose maximum value is the minimum value
of bin $sj$.

Figure~\ref{fig:binCountRangeSchematic} shows these calculations for
one `target' bin ($s=2, j=7$) and two
example species namely Moustache Sculpin and Snakeblenny (which are highlighted
as 127205 and 154675, respectively, in red in
Figure~\ref{fig:nSeaFung/nSeaFungBins1}). The resulting
minimum possible number (per hour) of individuals with body size larger than
those in bin $s=2, j=7$ is $E_{2,7,1} = 0.41$, and the maximum possible number
is $E_{2,7,2} = 2.71$, giving quite a range of possible values.

For the full data set for a particular year, similar calculations are done
for each combination of $s$ and $j$ (i.e. all bins for all species),
to give resulting ranges of counts $[E_{sj1}, E_{sj2}]$.
These are shown as the vertical span of each grey rectangle in
Figures~\ref{fig:ISD-1999.main} and \ref{fig:ISD-1986}-\ref{fig:ISD-2015} (one
figure for each year).
The horizontal span of the rectangle
depicts the range of that bin, i.e.~$[w_{sj}, w_{s,j+1}]$. The figures share a
common range for the x-axis.
%Note that
%we have not had to make any assumption about the distribution of individuals
%within bins.

The fitted curve, using the MLEbins method, is plotted as $(1 - F(x))n$, where
the total sample size is $n = \sum_{s=1}^S \sum_{j=1}^{J_s} d_{sj}$. The MLE for $b$
and the low and high values of the 95\% confidence intervals are also shown, but
the confidence intervals are fairly narrow and their
fits only just show up on the figures. Table~\ref{tab:annual.fit} gives,
for each year, the
estimated values of $\xmin$, $\xmax$, $n$, the MLE of $b$ and its 95\%
confidence interval, and the normalisation constant $C$.


% From nSeaFungMLEbinsRecommend-ISD.tex (need the \Sexpr{} evaluations here
Visually, the PLB model seems to fit fairly well (red curve passing through the
grey rectangles), at least
up to 100~g, for some years (in particular 1986, 1988, 1989, 1992, 1995,
1997, 1999, 2000, 2004, 2006, 2008, 2009, 2010 and 2014). Above 100~g the PLB
model generally overestimates the number of individuals. This could perhaps
be indicative of fishing pressure -- ideally we would have data concerning the
unfished community to see how well the PLB model fit such data. Developing a formal
statistical test of goodness-of-fit could be a focus of future work.

As well as no
significant trend of $b$ through time
(Figure~\ref{fig:nSeaFung/nSeaFungCompareTrends}),
Figures~\ref{fig:ISD-1986}-\ref{fig:ISD-2015} show no clear change in structure
through the 30~years of data. For example, there is not a disappearance of
larger fish through time.
For providing more rigorous advice to fisheries managers, assumptions that
could be investigated are the use of a minimum cut off of 4~g (based
on the minima of the body-mass bins), and the setting of $\xmin$ and $\xmax$ to
be the span of the data for each year (maybe only a smaller range is of interest,
and may provide a better fit of the PLB model).

% \subsection{Resulting plots and fits for each year of IBTS data\label{sec:resultfit}}

\clearpage

\thispagestyle{empty}
\vspace{-10mm}
% If edit this caption maybe also edit ISD-1999.main
\isdfig{1986}{Individual size distribution and MLEbins fit
            (red solid curve)
            with 95\% confidence intervals (red dashed curves)
            for IBTS data in 1986. The y-axis is linear in (a) and logarithmic
            in (b).
            For each bin, the horizontal green line shows
            the range of body sizes,
            with its value on the y-axis corresponding to the
            total number of individuals (per hour of trawling)
            in bins whose minima are $\geq$ the
            bin's minimum.
            The vertical span of each grey rectangle shows the
            possible range of the number of individuals with body mass
            $\geq$ the body mass of individuals in that bin (horizontal span
            is the same as for the green lines). The text in (a) gives the year,
            the MLE for the size-spectrum exponent $b$, and the sample size $n$.}

\newcounter{loop}
\newcommand{\loopMax}{2016}
\forloop{loop}{1987}{\value{loop} < \loopMax}
        {\isdfig{\arabic{loop}}{Individual size distribution and MLEbins fit
            with 95\% confidence intervals for IBTS data in \arabic{loop}.
            Details as in Figure~\ref{fig:ISD-1986}.
        }}

% From nSeaMLEbins-recommend.tex
% latex table generated in R 3.4.3 by xtable 1.8-2 package
% Thu May 17 15:32:32 2018
\begin{table}[ht]
\centering
\caption{Results for each year of the IBTS data from using the MLEbins method.
Both $\xmin$ and $\xmax$ are calculated as described earlier,
$n$ is the total number of individuals per hour,
`Low' and `High' give the 95\% confidence interval for the maximum likelihood
estimate (MLE) of $b$,
and $C$ is the normalisation constant from (\ref{PLB.C}) based on the MLE of $b$. }
% And $F(100) = \mbox{P}(X \leq 100)$ is the estimated proportion of individuals
%  that have body size $<100$~g.
%  Not including but can easily calculate.
\begin{tabular}{rrrrrrrr}
  \hline
Year & $\xmin$ & $\xmax$ & $n$ & Low & MLE of $b$ & High & $C$\\
  \hline
1986 & 4.05 & 28722.51 & 4799.29 & -1.50 & -1.49 & -1.47 & 0.97 \\
  1987 & 4.16 & 25974.17 & 6202.85 & -1.54 & -1.53 & -1.51 & 1.12 \\
  1988 & 4.06 & 29439.75 & 7711.72 & -1.62 & -1.60 & -1.59 & 1.41 \\
  1989 & 4.06 & 35210.99 & 7667.87 & -1.56 & -1.55 & -1.54 & 1.20 \\
  1990 & 4.05 & 34811.19 & 6399.75 & -1.53 & -1.52 & -1.50 & 1.07 \\
  1991 & 4.06 & 26412.52 & 6639.58 & -1.50 & -1.49 & -1.48 & 0.99 \\
  1992 & 4.05 & 25316.66 & 7973.32 & -1.54 & -1.53 & -1.52 & 1.12 \\
  1993 & 4.16 & 29439.75 & 8628.05 & -1.50 & -1.49 & -1.48 & 1.00 \\
  1994 & 4.16 & 22200.90 & 6282.84 & -1.53 & -1.52 & -1.50 & 1.10 \\
  1995 & 4.16 & 31818.64 & 9401.90 & -1.64 & -1.62 & -1.61 & 1.52 \\
  1996 & 4.05 & 36462.12 & 5607.19 & -1.44 & -1.42 & -1.41 & 0.78 \\
  1997 & 4.16 & 19360.99 & 9113.14 & -1.65 & -1.63 & -1.62 & 1.58 \\
  1998 & 4.06 & 22801.53 & 6927.12 & -1.57 & -1.55 & -1.54 & 1.22 \\
  1999 & 4.06 & 36462.12 & 7199.60 & -1.60 & -1.58 & -1.57 & 1.33 \\
  2000 & 4.06 & 22801.53 & 10816.80 & -1.61 & -1.60 & -1.59 & 1.41 \\
  2001 & 4.05 & 18824.87 & 9058.45 & -1.46 & -1.45 & -1.44 & 0.85 \\
  2002 & 4.06 & 20464.76 & 6766.12 & -1.44 & -1.42 & -1.41 & 0.79 \\
  2003 & 4.05 & 21611.31 & 5614.19 & -1.45 & -1.44 & -1.42 & 0.82 \\
  2004 & 4.05 & 30911.23 & 5445.26 & -1.52 & -1.50 & -1.49 & 1.02 \\
  2005 & 4.05 & 19360.99 & 4189.79 & -1.39 & -1.38 & -1.36 & 0.67 \\
  2006 & 4.05 & 21032.63 & 7864.87 & -1.63 & -1.62 & -1.60 & 1.48 \\
  2007 & 4.06 & 18299.10 & 5278.55 & -1.52 & -1.50 & -1.49 & 1.03 \\
  2008 & 4.06 & 28722.51 & 6516.37 & -1.57 & -1.56 & -1.54 & 1.23 \\
  2009 & 4.05 & 24036.35 & 5628.83 & -1.59 & -1.58 & -1.56 & 1.30 \\
  2010 & 4.06 & 23293.68 & 8049.43 & -1.54 & -1.53 & -1.52 & 1.13 \\
  2011 & 4.06 & 26723.46 & 11259.91 & -1.49 & -1.48 & -1.47 & 0.94 \\
  2012 & 4.06 & 14899.37 & 7670.79 & -1.53 & -1.51 & -1.50 & 1.07 \\
  2013 & 4.08 & 25974.17 & 5736.94 & -1.57 & -1.56 & -1.54 & 1.23 \\
  2014 & 4.06 & 30025.12 & 8122.67 & -1.63 & -1.62 & -1.61 & 1.48 \\
  2015 & 4.08 & 18362.51 & 9072.14 & -1.76 & -1.75 & -1.73 & 2.14 \\
   \hline
\end{tabular}
\label{tab:annual.fit}
\end{table}

\clearpage

\subsection{Further results from fitting simulated data\label{sec:furtherSim}}

Table~\ref{tab:rep} gives the summary statistics for the simulation results
shown in Figure~\fitMLEhists. The mean and median estimates of $b$
(that has true value of -2) for the MLEbin method are all -1.99 or -2.00. But
for the MLEmid method the estimates are less accurate and depend upon the bin
width.
% Copying from the output from running
%  fitting-size-spectra/codeBinStructure/binFitting/fitMLEbinConf.r (i.e. not
%  done as xtable).
\begin{table}[b]
\caption{\label{tab:rep} Summary statistics for the 10,000 simulations of 1,000
  samples from (\ref{PLB}), corresponding to the histograms in
  Figure~\fitMLEhists. Each sample is binned using each of four binning types:
  `Linear 1', `Linear 5' and `Linear 10' for bins of constant bin widths 1, 5,
  and 10, respectively, and `2k' for bin widths that double in size. Each binned
  sample is then fit using the MLEmid and MLEbin methods, where the MLEmid
  method uses the midpoint of each bin and the MLEbin method (shaded rows)
  more explicitly
  accounts for the binning in the likelihood function. Statistics relate to the
  resulting 10,000 estimates of $b$, with the final column giving the percentage
  of simulations for which the estimate is below the true value of $b = -2$.}
\rowcolors{3}{}{lightgray}
\begin{tabular}{llrrrrr}
\hline Binning type & Method & 5\% quantile & Median & Mean & 95\% quantile &
Percentage\\ & & & & & & below true\\ \hline Linear 1 & MLEmid & -1.98 & -1.94 &
-1.94 & -1.89 & 1\\ & MLEbin & -2.05 & -1.99 & -1.99 & -1.94 & 43 \\ Linear 5 &
MLEmid & -1.63 & -1.61 & -1.60 & -1.57 & 0 \\ & MLEbin & -2.06 & -1.99 & -1.99 &
-1.94 & 42 \\ Linear 10 & MLEmid & -1.43 & -1.41 & -1.40 & -1.35 & 0 \\ & MLEbin
& -2.06 & -1.99 & -1.99 & -1.93 & 42 \\ 2k & MLEmid & -1.95 & -1.90 & -1.90 &
-1.86 & 0 \\ & MLEbin & -2.05 & -2.00 & -2.00 & -1.94 & 47\\ \hline
\end{tabular}
\end{table}

We repeat those simulations but with $\xmin = 16$ rather than $\xmin=1$ to
test the sensitivity of the methods.
Figures~\fitMLEhistsXminSixteen~and \fitMLEconfsXminSixteen~show the equivalent
results to Figures~\fitMLEhists~and \fitMLEconfs, and Table~\ref{tab:repXmin16}
gives the summary statistics. The MLEmid method behaves better than it did
for $\xmin=1$, but not as well as the MLEbin methods that still gives
consistently good results.

We also simulate a community with $\xmin=1$, but with the observed data
consisting of all body masses $\geq c$, where $c$
is a cutoff value. In practice, $c$ could be known from the sampling protocol or
assigned a value based on previous experience. Sufficient random body masses
were sampled to obtain a sample size of $n=1,000$ for the observed data, and we
set $c=16$.

The results (Table~\ref{tab:repCutOff16} and Figures~\fitMLEhistsCutOffSixteen\
and~\fitMLEconfsCutOffSixteen) are essentially identical to those for $\xmin=16$.
This emphasises the scale-free nature of the power-law
distribution, whereby the cutoff value does not affect the estimate of $b$.
Whereas for a scale-dependent distribution, such as the lognormal,
the estimated parameters would depend on the cutoff value.

\begin{table}
\caption{\label{tab:repXmin16} As for Table \ref{tab:rep} but with $\xmin = 16$,
  corresponding to the histograms in Figure~\fitMLEhistsXminSixteen .}
\rowcolors{3}{}{lightgray}
\begin{tabular}{llrrrrr}
\hline Binning type & Method & 5\% quantile & Median & Mean & 95\% quantile &
Percentage\\ & & & & & & below true\\ \hline Linear 1 & MLEmid & -2.06 & -1.99 &
-2.00 & -1.94 & 44 \\ & MLEbin & -2.06 & -1.99 & -2.00 & -1.94 & 45 \\ Linear 5
& MLEmid & -2.05 & -1.99 & -1.99 & -1.93 & 36 \\ & MLEbin & -2.06 & -1.99 &
-2.00 & -1.94 & 45 \\ Linear 10 & MLEmid & -2.02 & -1.96 & -1.96 & -1.90 & 14
\\ & MLEbin & -2.06 & -2.00 & -2.00 & -1.94 & 45 \\ 2k & MLEmid & -1.93 & -1.87
& -1.87 & -1.82 & 0 \\ & MLEbin & -2.07 & -2.00 & -2.00 & -1.94 & 51 \\ \hline
\end{tabular}
\end{table}


\begin{table}
\caption{\label{tab:repCutOff16} As for Table \ref{tab:rep} with $\xmin = 1$,
  but only sampling data above the cutoff value (i.e.~${\geq}c=16$),
  corresponding to the histograms in Figure~\fitMLEhistsCutOffSixteen . Each
  simulated data set had a sample size of 1,000 above the cutoff value. Results
  are identical to those in Table~\ref{tab:repXmin16} except for a few
  differences of 0.01 in the statistical values and 1\% in the final column.}
\rowcolors{3}{}{lightgray}
\begin{tabular}{llrrrrr}
\hline Binning type & Method & 5\% quantile & Median & Mean & 95\% quantile &
Percentage\\ & & & & & & below true\\ \hline
% source("fitMLEbinConfCutOff16.r")
% Overall range of MLE values:
% -2.172917 -1.745742
% Overall range of confidence intervals:
% -2.256917 -1.682742
Linear 1 & MLEmid & -2.06 & -1.99 & -2.00 & -1.93 & 45 \\ & MLEbin & -2.06 &
-1.99 & -2.00 & -1.94 & 45 \\ Linear 5 & MLEmid & -2.05 & -1.99 & -1.99 & -1.93
& 36 \\ & MLEbin & -2.06 & -2.00 & -2.00 & -1.94 & 45 \\ Linear 10 & MLEmid &
-2.02 & -1.96 & -1.96 & -1.90 & 13 \\ & MLEbin & -2.06 & -2.00 & -2.00 & -1.93 &
46 \\ 2k & MLEmid & -1.93 & -1.87 & -1.87 & -1.82 & 0 \\ & MLEbin & -2.07 &
-2.00 & -2.00 & -1.94 & 51 \\
% Binning type & Method & Shortest CI & Widest CI\\ \\hline
% Linear 1  &  MLEmid  &  0.134  &  0.162 \\
%  &  MLEbin  &  0.134  &  0.162 \\
% Linear 5  &  MLEmid  &  0.134  &  0.161 \\
%  &  MLEbin  &  0.135  &  0.162 \\
% Linear 10  &  MLEmid  &  0.133  &  0.158 \\
%  &  MLEbin  &  0.136  &  0.165 \\
% 2k  &  MLEmid  &  0.127  &  0.146 \\
%  &  MLEbin  &  0.136  &  0.165 \\
\hline
\end{tabular}
\end{table}

\clearpage

\thispagestyle{empty}

% \fitMLEhistsXminSixteen
\onefig{binFitting/xmin16/fitMLEhistsXmin16}{As for
  Figure~\fitMLEhists~but with $\xmin = 16$.}

\clearpage
\thispagestyle{empty}

% \fitMLEconfsXminSixteen
\onefigshift{binFitting/xmin16/fitMLEconfsXmin16}{As for
  Figure~\fitMLEconfs~but with $\xmin = 16$.}

\clearpage

\thispagestyle{empty}

% \fitMLEhistsCutOffSixteen
\onefig{binFitting/xCutOff16/fitMLEhistsCutOff16}{As for
  Figure~\fitMLEhists~but only observing data above a cutoff of 16. There are
  only very minor differences to Figure~\fitMLEhistsXminSixteen, as confirmed in
  Table~\ref{tab:repCutOff16}.}
% Histograms of estimated exponent $b$ for 10,000 simulated data sets, each of which contains 1,000 independent random numbers drawn from a bounded power-law distribution with $b=-2$, $\xmin=16$ and $\xmax=1,000$. Panels (a)-(f) use linear bins of constant bin width 1, 5 or 10 (binning type `Linear 1' etc.), and panels (g)-(h) use bin widths that double in size (binning type `2k'). Each binned data set is fitted using either the MLEmid method (a, c, e, g) or the MLEbin method (b, d, f, h). Vertical red lines indicate the known value of $b=-2$.}

\clearpage
\thispagestyle{empty}

% \fitMLEconfsCutOffSixteen
\onefigshift{binFitting/xCutOff16/fitMLEconfsCutOff16}{As for
  Figure~\fitMLEconfs~but only observing data $\geq$ the cutoff value of
  16. There are only very minor differences to Figure~\fitMLEconfsXminSixteen,
  with five of the observed coverage values changing by 1\%.}

\clearpage

\subsection{Histograms using different binning types\label{sec:hists}}

Figures~\ref{fig:binFitting/exampleHists/histBinType1}-\ref{fig:binFitting/exampleHists/histBinType2kdensity}
show histograms of the same set of 1,000 random numbers
sampled from the PLB distribution (with $b=-2$, $\xmin=1$
and $\xmax=1,000$), plotted using the different binning types
described in the main text.
These illustrate the highly skewed nature of the PLB distribution, and how
the first bin can contain most of the counts.

\onefig{binFitting/exampleHists/histBinType1}{Histogram of the 1,000 random numbers with bin
  widths of 1.}

\onefig{binFitting/exampleHists/histBinType5}{Histogram of the 1,000 random numbers with bin
  widths of 5.}

\onefig{binFitting/exampleHists/histBinType10}{Histogram of the 1,000 random numbers with bin
  widths of 10.}

\onefig{binFitting/exampleHists/histBinType2k}{Histogram of the 1,000 random numbers with bin
  widths that double in size. Note that heights of bars represent the counts in
  each bin, and that because of the non-constant bin widths the areas are not
  proportional to the counts -- see
  Figure~\ref{fig:binFitting/exampleHists/histBinType2kdensity}.}

\onefig{binFitting/exampleHists/histBinType2kdensity}{Density plot of the 1,000 random
  numbers with bin widths that double in size, showing density rather than
  frequency.}

% \onefig{binFitting/exampleHists/histBinType50}{Histogram of the 1,000 random numbers with bin widths of 50.}

\renewcommand{\bibsection}{\subsection{Literature cited only in Supplementary Material}}

\begin{thebibliography}{5}

\bibitem[Burnham and Anderson(2002)]{ba02}
Burnham, K.~P. and Anderson, D.~R. (2002).
\newblock {\em Model Selection and Multimodel Inference: A Practical
  Information-Theoretic Approach}.
\newblock 2nd ed., Springer, New York.

\bibitem[Edwards et~al.~(2007)]{edwa07}
Edwards, A.~M., Phillips, R.~A., Watkins, N.~W., Freeman, M.~P., Murphy, E.~J.,
  Afanasyev, V., Buldyrev, S.~V., da~Luz, M. G.~E., Raposo, E.~P., Stanley,
  H.~E., \& Viswanathan, G.~M. (2007).
\newblock Revisiting {L}\'evy flight search patterns of wandering albatrosses,
  bumblebees and deer.
\newblock {\em Nature}, {\bf 449}, 1044--1048.

\bibitem[Edwards et~al.~(2012)]{efbj12}
Edwards, A.~M., Freeman, M.~P., Breed, G.~A., \& Jonsen, I.~D. (2012).
\newblock Incorrect likelihood methods were used to infer scaling laws of
  marine predator search behaviour.
\newblock {\em PLoS ONE}, {\bf 7}, e45174.

\end{thebibliography}{}

\section{Summary of code\label{sec:code}}

For full reproducibility, all R (and Sweave) code is freely available at

\noindent {\protect{\tt
    https://github.com/andrew-edwards/fitting-size-spectra/tree/manuscript2}}. Note
that this is the code repository associated with \citet{\MEE},
but with new code in the branch
{\tt manuscript2} that will be merged into the the {\tt master} branch in due
course. The code for this manuscript relies on some of the functions from the
\citet{\MEE} code. The full repository is also provided in the submitted file
{\tt fit-bin-submit.zip}. To get started, please see the file {\tt codeBinStructure/readMeBinStructure.txt}.





%TC:endignore

\end{document}
